\documentclass[11pt]{article}

\usepackage[margin=1in]{geometry}
\usepackage{amsmath,amssymb,amsthm,mathtools}
\usepackage{booktabs}
\usepackage{array}
\usepackage{enumitem}
\usepackage[colorlinks=true,linkcolor=blue,citecolor=blue]{hyperref}

% ---- notation contract (used uniformly throughout) ----
\newcommand{\Z}{\mathbb{Z}}
\newcommand{\F}[1]{\mathbb{F}_{#1}}
\newcommand{\wt}[1]{\lvert #1\rvert}
\DeclareMathOperator{\Ann}{Ann}
\DeclareMathOperator{\img}{im}
\DeclareMathOperator{\supp}{supp}
\DeclareMathOperator{\Stab}{Stab}
\newcommand{\aug}{\varepsilon}
\newcommand{\hA}{\hat{A}}
\newcommand{\hB}{\hat{B}}
\newcommand{\dcut}{\partial_2^{\mathrm{c}}}     % seam-crossing part of the cover boundary
\newcommand{\dncut}{\partial_2^{\mathrm{nc}}}   % non-crossing part
\newcommand{\cmap}{\Delta}                       % Smith connecting map
\newcommand{\slab}{\ell}                          % slot-labeling vector (kappa of a multiplier)
\newcommand{\sx}{s_x}
\newcommand{\sy}{s_y}
\newcommand{\tx}{t_x}
\newcommand{\ty}{t_y}

\theoremstyle{plain}
\newtheorem*{thmA}{Theorem A}
\newtheorem*{thmB}{Theorem B}
\newtheorem*{thmC}{Theorem C}
\newtheorem*{thmD}{Theorem D}
\newtheorem*{thmR}{Theorem (R)}
\newtheorem{lemma}{Lemma}[section]
\newtheorem{proposition}[lemma]{Proposition}
\newtheorem{corollary}[lemma]{Corollary}
\theoremstyle{definition}
\newtheorem{definition}[lemma]{Definition}
\theoremstyle{remark}
\newtheorem*{remark}{Remark}

\title{A self-contained analytic proof that the gross code\\ has distance $12$}
\author{}
\date{}

\begin{document}
\maketitle

\begin{abstract}
The \emph{gross code} is the $[[144,12,12]]$ bivariate-bicycle (BB) CSS code of
Bravyi, Cross, Gambetta, Maslov, Rall and Yoder, defined over the group
$\Z_{12}\times\Z_6$ by the polynomials $A=x^3+y+y^2$ and $B=y^3+x+x^2$. We give a
complete, self-contained proof that its minimum distance is exactly $12$. The argument is
analytic: it is organized around the realization of the gross code as the free $\Z_2$
double cover, in the $x$-direction, of the $[[72,12,6]]$ \emph{base} code over
$\Z_6\times\Z_6$ with the same $A,B$. Every finite case analysis is small enough to be
surveyed by hand. The two ingredients are a Chinese-remainder (CRT) ``layer frame'' that
diagonalizes the relevant convolution operators over $\F4$, and a cover-transfer
dichotomy that splits each logical class into a \emph{dangerous} part (controlled by a
factor-two lemma against the base code's stabilizer structure) and a \emph{safe} part
(controlled by a confined-frame weight floor). Along the way we prove $d(\mathrm{base})=6$
and $d(\mathrm{gross})\ge 6$, the latter already tripling the best general lower bound
available from the degeneracy-divided product-code estimate.
\end{abstract}

\section{Results}\label{sec:results}

Fix the finite abelian group and the two polynomials of the family. For the
\emph{base} code take $G=\Z_6\times\Z_6$ and for the \emph{gross} code take
$G=\Z_{12}\times\Z_6$; in both cases
\[
  A=x^3+y+y^2,\qquad B=y^3+x+x^2\qquad\in\ \F2[G],
\]
where a polynomial is identified with the indicator function of its support (the monomial
$x^a y^b$ is the point $(a,b)\in G$). We prove four results, referred to throughout as
Theorems~A--D.

\begin{thmA}[no small cycles]
The base code $[[72,12,6]]$ has no nonzero $X$-type or $Z$-type $1$-cycle of weight
$\le 5$. In particular $d(\mathrm{base})\ge 6$, and the minimum weights $\mu_Z,\mu_X$ of a
nonzero $Z$- resp.\ $X$-stabilizer are $\ge 6$.
\end{thmA}

\noindent\textbf{Corollary A$'$.} $d(\mathrm{base})=6$: the weight-$6$ element
$z^\ast=1+y+y^2+y^5+x^3+x^3y^4\in\Ann(A)$ furnishes the $Z$-logical $(z^\ast,0)$
(\S\ref{sec:thmA}).

\begin{thmB}[the cover floor]
The gross code $[[144,12,12]]$ --- the free $\Z_2$ double cover of the base in the
$x$-direction --- satisfies $d(\mathrm{gross})\ge 6$.
\end{thmB}

\begin{thmC}[the dangerous sector, tight]
Every nontrivial $Z$-logical of the gross code whose class lies in the kernel of the
projection $p_\ast\colon H_1(\mathrm{gross})\to H_1(\mathrm{base})$ has weight
$\ge 12=2\,d(\mathrm{base})$, and weight $12$ is attained.
\end{thmC}

\begin{thmD}[the distance]
$d(\mathrm{gross})=12$.
\end{thmD}

The distance values $6$, $6$, $12$ themselves are not new --- they are accessible to
exact integer-programming solvers. What the argument supplies is a structural
\emph{mechanism}: a cover-transfer dichotomy through which the base code's geometry
controls the gross code's distance, by routes that do not require any exhaustive search
over the $2^{144}$ logical operators.

\section{Setup and conventions}\label{sec:setup}

A bivariate-bicycle code over a finite abelian group $G$, with $A,B\in\F2[G]$, has qubit
set $L\sqcup R$, each block a copy of $\F2[G]$, and check matrices
$H_X=(M_A\mid M_B)$, $H_Z=(M_B^{\mathsf T}\mid M_A^{\mathsf T})$, where $M_P$ is the
operator of multiplication (convolution) by $P$ in $\F2[G]$. We use the chain complex
\[
  C_2\ \xrightarrow{\ \partial_2\ }\ C_1\ \xrightarrow{\ \partial_1\ }\ C_0,
  \qquad \partial_1=H_X,\quad \partial_2=H_Z^{\mathsf T},
\]
so that for $z\in\F2[G]$,
\[
  \partial_2 z=(B\cdot z,\ A\cdot z)\qquad\text{(left block $B\cdot z$, right block $A\cdot z$)} .
\]
Here $P\cdot f$ denotes the convolution product $(P\cdot f)(g)=\sum_{h}P(h)f(g-h)$. The
$Z$-stabilizers are $\img\partial_2$, the $Z$-logical operators are
$\ker\partial_1\smallsetminus\img\partial_2$, and the $Z$-distance is the minimum weight
over the nonzero classes of $H_1=\ker\partial_1/\img\partial_2$:
\[
  d_Z=\min\{\,\wt{u} : u\in\ker\partial_1,\ u\notin\img\partial_2\,\}.
\]
A $1$-chain is written $u=(u_L,u_R)$ with $u_L,u_R\in\F2[G]$; the cycle condition
$\partial_1 u=0$ reads
\[
  A\cdot u_L=B\cdot u_R\ (=:\sigma),\qquad\text{``the matched syndrome.''}
\]
We write $\wt{x}$ for Hamming weight. The \emph{augmentation}
$\aug\colon\F2[G]\to\F2$ (sum of coefficients) is a ring homomorphism with
$\aug(A)=\aug(B)=1$, and weight is congruent to augmentation modulo $2$; hence
\begin{equation}\label{eq:PAR}
  \wt{A\cdot f}\equiv\wt{f}\quad\text{and}\quad \wt{B\cdot f}\equiv\wt{f}\pmod 2.\tag{PAR}
\end{equation}

\paragraph{Difference sets.}
Write $\mathrm d A=\{g-h: g\ne h\in\supp A\}$ and similarly $\mathrm d B$. In coordinates
$(x,y)$,
\[
  \mathrm d A=\{(0,\pm1),(3,\pm1),(3,\pm2)\},\qquad
  \mathrm d B=\{(\pm1,0),(\pm1,3),(\pm2,3)\}.
\]
Both are multiplicity-free (each element arises from exactly one ordered pair), they are
disjoint, and they are disjoint \emph{in each coordinate} separately:
$x(\mathrm dA)\subseteq\{0,3\}$, $x(\mathrm dB)\subseteq\{1,2,4,5\}$,
$y(\mathrm dA)\subseteq\{1,2,4,5\}$, $y(\mathrm dB)\subseteq\{0,3\}$. Consequently any two
columns of $M_A$ (or any two of $M_B$) intersect in at most one cell (``$\mathrm{ov}\le1$''),
with equality exactly on $\mathrm dA$ (resp.\ $\mathrm dB$).

\section{The \texorpdfstring{$X$--$Z$}{X-Z} inversion duality}\label{sec:duality}

\begin{lemma}[duality]\label{lem:duality}
For any BB code (any abelian $G$, any $A,B$), the map
$\Phi(w_L,w_R)=(\iota(w_R),\iota(w_L))$, where $\iota$ is the antipode $g\mapsto -g$
extended $\F2$-linearly, is a weight-preserving involution carrying $\ker H_Z$ onto
$\ker H_X$ and the $X$-stabilizer space $\img H_X^{\mathsf T}$ onto the $Z$-stabilizer
space $\img H_Z^{\mathsf T}$. Hence $d_X=d_Z$.
\end{lemma}

\begin{proof}
Since $G$ is abelian, $\iota$ is an algebra automorphism of $\F2[G]$, and
$M_P^{\mathsf T}=M_{\iota(P)}$. A vector $(w_L,w_R)\in\ker H_Z$ satisfies
$\iota(B)\,w_L+\iota(A)\,w_R=0$; applying $\iota$ gives $B\cdot\iota(w_L)+A\cdot\iota(w_R)=0$,
i.e.\ $\Phi(w)\in\ker H_X$. The row of $H_X$ at check $g$ is $(\iota(A)\,\delta_g,\iota(B)\,\delta_g)$,
and $\Phi$ sends it to $(B\cdot\delta_{-g},A\cdot\delta_{-g})=\partial_2\delta_{-g}$; thus
spans map onto spans. Finally $\Phi^2=\mathrm{id}$ and $\iota$ preserves weight.
\end{proof}

This applies to both the base and the gross codes (inversion needs no symmetry between the
$x$- and $y$-directions). \emph{All statements below are therefore proved on the $Z$ side
only}, the $X$ side following by Lemma~\ref{lem:duality}.
\section{The CRT layer frame for the base}\label{sec:crtframe}

All structure in the dangerous sector flows from a single algebraic
observation: the group algebra $\F2[\Z_6^2]$ splits, by the Chinese remainder
theorem, into a layer part $\F2[\Z_2^2]$ and four $\F4$-components on which the
multipliers $A,B$ act as small, explicit $\F4[\Z_2^2]$-elements. We set up that
frame here and extract the two structural facts the cover argument consumes:
the \emph{Engine Lemma}, governing the radical multipliers on each component,
and the \emph{layer dictionary}, governing minimum weights on the
$\Z_3^2$-fibers.

\subsection{The splitting}

Write $\Z_6=\Z_2\times\Z_3$ in each coordinate, with $x=\sx\,\tx$ where
$\sx=x^3$ generates the $2$-part and $\tx=x^4$ generates the $3$-part, and
likewise $y=\sy\,\ty$. A cell of $\Z_6^2$ is then a pair $(s,t)$ with
$s=(\sx,\sy)\in\Z_2^2$ (the \emph{layer}) and $t=(\tx,\ty)\in\Z_3^2$. Over
$\F4$ the eight characters of the $3$-part $\Z_3^2$ split into five Frobenius
orbits, and correspondingly
\begin{equation}\label{eq:crt-split}
  R:=\F2[\Z_6^2]\ \cong\ R_0\times R_1\times R_2\times R_3\times R_4,
  \qquad R_0=\F2[\Z_2^2],\quad R_j=\F4[\Z_2^2]\ (j=1,\dots,4),
\end{equation}
where $R_0$ is the trivial-orbit (layer) part and each $R_j$ carries the
$\F4[\Z_2^2]$-structure of one nontrivial Frobenius orbit of $3$-part
characters. As functions of $t=(\tx,\ty)$ the component characters are
\[
  \psi_1=\omega^{\ty},\quad
  \psi_2=\omega^{\tx},\quad
  \psi_3=\omega^{\tx+\ty},\quad
  \psi_4=\omega^{\tx+2\ty},
\]
where $\omega$ is a fixed primitive cube root of unity in $\F4$. These satisfy
the multiplicative relations
\[
  \psi_3=\psi_1\psi_2,\qquad \psi_4=\psi_1\psi_3,
\]
and have kernels (as subgroups of $\Z_3^2$)
\[
  \ker\psi_1=\langle(1,0)\rangle,\quad
  \ker\psi_2=\langle(0,1)\rangle,\quad
  \ker\psi_3=\langle(1,2)\rangle,\quad
  \ker\psi_4=\langle(1,1)\rangle.
\]

\subsection{The component multipliers}

On the layer algebra $\F2[\Z_2^2]$ put
\[
  u:=1+\sx,\qquad v:=1+\sy,\qquad u^2=v^2=0,
\]
so $u,v$ are the two square-zero generators of the radical of $\F2[\Z_2^2]$,
and $uv=1+\sx+\sy+\sx\sy$ has the all-ones coefficient vector over the four
layers. For each component $j$ let
$(\xi_j,\eta_j)=\bigl(\psi_j(\tx\text{-unit}),\psi_j(\ty\text{-unit})\bigr)\in\F4^2$
record the values of $\psi_j$ on the two generators of $\Z_3^2$. Pushing $A$ and
$B$ through the splitting~\eqref{eq:crt-split} and writing $\xi=\xi_j$,
$\eta=\eta_j$, the component images of the multipliers are
\begin{equation}\label{eq:crt-mult}
  \hA_j=(1+\eta+\eta^2)+u+\eta v,\qquad
  \hB_j=(1+\xi+\xi^2)+v+\xi u\qquad\in\ \F4[\Z_2^2].
\end{equation}
Since $1+\zeta+\zeta^2=0$ for $\zeta\in\{\omega,\omega^2\}$ and equals $1$ for
$\zeta=1$, the constant terms vanish exactly when the corresponding generator is
nontrivial on the orbit. Tabulating over the five orbits:

\begin{center}
\begin{tabular}{cccc}
\toprule
$j$ & $(\xi_j,\eta_j)$ & $\hA_j$ & $\hB_j$ \\
\midrule
$0$ & $(1,1)$        & $1+u+v$ \ (unit)   & $1+u+v$ \ (unit) \\
$1$ & $(1,\omega)$   & $u+\omega v$ \ (radical) & $1+u+v$ \ (unit) \\
$2$ & $(\omega,1)$   & $1+u+v$ \ (unit)   & $\omega u+v$ \ (radical) \\
$3$ & $(\omega,\omega)$  & $u+\omega v$   & $\omega u+v$ \\
$4$ & $(\omega,\omega^2)$ & $u+\omega^2 v$ & $\omega u+v=\omega\cdot\hA_4$ \\
\bottomrule
\end{tabular}
\end{center}

\noindent
The component transform diagonalizing convolution is read off the splitting as
follows. For $f\in\F2[\Z_6^2]$ and a layer $s\in\Z_2^2$ write $f_s\in\F2[\Z_3^2]$
for the restriction of $f$ to the fiber over $s$. The transform of $f$ on
component $j$ is the vector $V_j(f)\in\F4^{\Z_2^2}$,
\[
  V_j(f)[s]=\sum_{t\in\supp f_s}\psi_j(t),
\]
and $V_0(f)[s]$ is the $\F2$-parity of the layer $f_s$. Convolution becomes
pointwise multiplication: with $\widehat{z}_j:=V_j(z)$,
\begin{equation}\label{eq:crt-mtable}
  V_j(A\cdot z)=\hA_j\cdot\widehat{z}_j,\qquad
  V_j(B\cdot z)=\hB_j\cdot\widehat{z}_j
  \qquad\text{(products in }\F4[\Z_2^2]\text{)}.
\end{equation}
This is the layer frame: each convolution operator on $\F2[\Z_6^2]$ is, component
by component, multiplication by one of the five $\F4[\Z_2^2]$-elements above.

\subsection{The Engine Lemma}

The two diagonal components $R_3,R_4$ and the two single-radical components
$R_1,R_2$ all involve multiplication by one of the six \emph{radical}
multipliers
\[
  u+\omega v,\quad u+\omega^2 v,\quad \omega u+v,\quad
  \omega^2 u+v,\quad u+v,\quad \text{(and the } \hA,\hB \text{ entries above)},
\]
i.e.\ a nonzero element of the form $D=\alpha u+\beta v$ with $\alpha,\beta\in\F4$
not both zero and $D$ not a unit. The next lemma describes the annihilator and
support behavior of any such $D$; it is the workhorse (``engine'') behind every
support count in the dangerous sector.

\begin{lemma}[engine]\label{lem:engine}
Let $D\in\F4[\Z_2^2]$ be any of the six radical multipliers appearing
in~\eqref{eq:crt-mult}, i.e.\ $D=\alpha u+\beta v$ with $(\alpha,\beta)\ne(0,0)$
and $D$ a zero divisor.
\begin{enumerate}
\item[\textup{(i)}] \emph{Annihilator.} $\Ann(D)=(D)=\{\alpha' D+\beta'\,uv:
  \alpha',\beta'\in\F4\}$, a $2$-dimensional ideal. In particular $D^2=0$ and
  $D\cdot uv=0$.
\item[\textup{(ii)}] \emph{Value vector and support dichotomy.} Evaluated on the
  four layers $(1,\sx,\sy,\sx\sy)$, the multiplier $D$ has value vector of the
  shape $(1+\eta,1,\eta,0)$ for a suitable $\eta\in\{\omega,\omega^2\}$: three
  pairwise-distinct nonzero values and exactly one zero. Consequently a nonzero
  ideal element $g=\alpha' D+\beta'\,uv\in(D)$ has support either
  \begin{itemize}
    \item \textbf{full} ($\alpha'=0$): $g=\beta'\,uv$ is the constant vector
      $\beta'\,\vec{1}$ on all four layers, since $uv$ has the all-ones
      coefficient vector; or
    \item a \textbf{co-point} ($\alpha'\ne0$): $g$ vanishes at the single layer
      $s_\star$ where $\alpha' D=\beta'$ and is nonzero on the other three, with
      value vector $\alpha'\,C(s_\star)$, where $C(s_\star):=D+D[s_\star]\,\vec{1}$,
      rigid up to the overall scalar $\alpha'$.
  \end{itemize}
\item[\textup{(iii)}] \emph{$C$-table.} Writing $\eta_j\in\{\omega,\omega^2\}$ for
  the component value and using $\eta^2=1+\eta$, the basepoint co-point profile is
  \[
    C_j([1])=(0,\ \eta_j,\ 1,\ \eta_j^2)\qquad\text{over }(1,\sx,\sy,\sx\sy),
  \]
  and in general $C_j(s_\star)[s]=\eta_j^{\,e(s_\star,s)}$ with an exponent
  $e(s_\star,s)$ that is \emph{independent of $j$}. Hence every cross-layer ratio
  of $C$-values is a power of $\eta_j$ with a $j$-independent exponent. Because
  $\psi_4=\psi_1\psi_3$ and $\eta_4=\eta_1\eta_3$, any ratio system
  ``$\psi_j(\tau)=$ ($C$-ratio)$_j$'' for $j\in\{1,3,4\}$ is automatically
  consistent and pins $\tau$ to a multiple of $(0,1)$ (the characters $\psi_3$,
  $\psi_4$ separate $\Z_3^2$).
\end{enumerate}
\end{lemma}

\begin{proof}
\emph{(i)} Each radical multiplier satisfies $D^2=\alpha^2 u^2+\beta^2
v^2=0$ since $u^2=v^2=0$ and the cross term $2\alpha\beta\,uv$ vanishes in
characteristic $2$. Also $D\cdot uv=\alpha u\cdot uv+\beta v\cdot uv=0$ because
$u^2v=uv^2=0$. Thus $D$ and $uv$ both lie in $\Ann(D)$, and they are
$\F4$-independent (one is supported in the radical-degree-$1$ part, the other is
$uv$). The radical $\mathfrak m=(u,v)$ of $\F4[\Z_2^2]$ has $\F4$-dimension $3$
(basis $u,v,uv$) and $\mathfrak m^2=(uv)$ is $1$-dimensional; for a degree-$1$
radical element $D$ the principal ideal $(D)=\F4 D+\F4\,uv$ is exactly
$2$-dimensional. A dimension count then gives $\Ann(D)=(D)$: multiplication by
$D$ kills $(D)$ by the relations just shown, while its image is
$\operatorname{span}_{\F4}\{D,uv\}$ --- from $D\cdot 1=D$ together with
$D\cdot u=\beta\,uv$ and $D\cdot v=\alpha\,uv$ --- of $\F4$-dimension $2$. Rank--nullity
on the $4$-dimensional algebra $\F4[\Z_2^2]$ then forces $\dim_{\F4}\Ann(D)=4-2=2$,
so $\Ann(D)=(D)$.

\emph{(ii)} Evaluate $D=\alpha u+\beta v=\alpha(1+\sx)+\beta(1+\sy)$ on the four
layers using the layer characters $s\mapsto(\pm1,\pm1)$. On the trivial layer
$s=1$ one gets $\alpha\cdot 0+\beta\cdot 0$ in the $\sx,\sy$ slots after
subtracting the constant, producing the value $1+\eta$ in the normalization where
$\eta=\beta/\alpha$ (or its inverse); on the $\sx$-flip layer $D$ takes value
$1$, on the $\sy$-flip layer it takes $\eta$, and on $\sx\sy$ it takes $0$. These
are three distinct nonzero entries of $\F4^\times=\{1,\omega,\omega^2\}$ and one
zero. Now let $g=\alpha'D+\beta'uv\in(D)$ be nonzero. If $\alpha'=0$ then
$g=\beta'uv$, whose coefficient vector over the layers is $\beta'$ times the
all-ones vector, hence full support and constant value. If $\alpha'\ne0$, then
$g$ is the layer-function $\alpha'D+\beta'\vec 1$ (since $uv\mapsto\vec 1$ on
layers); it vanishes exactly where $\alpha'D=\beta'$, and by the distinctness of the
nonzero values established above there is exactly one such layer $s_\star$, while the remaining three layers
carry the nonzero values $\alpha'\bigl(D[s]-D[s_\star]\bigr)$. Factoring out
$\alpha'$ leaves the vector $C(s_\star)=D+D[s_\star]\vec 1$, which is determined
by $s_\star$ and rigid up to scalar.

\emph{(iii)} Specializing the value vector of (ii) at the basepoint $s_\star=[1]$
gives $C_j([1])=(0,\eta_j,1,\eta_j^2)$ once we use $\eta_j^2=1+\eta_j$ for
$\eta_j\in\{\omega,\omega^2\}$. Translating $s_\star$ by a layer simply permutes
the four entries by the corresponding $\Z_2^2$-shift, replacing each value by
$\eta_j$ raised to a shifted exponent; the bookkeeping of which power occurs at
which layer depends only on the relative layer displacement $e(s_\star,s)$, not on
which component $j$ we are in. Thus $C_j(s_\star)[s]=\eta_j^{\,e(s_\star,s)}$ with
$e$ independent of $j$, and every cross-layer ratio
$C_j(s_\star)[s]/C_j(s_\star)[s']$ is $\eta_j$ to a $j$-independent power. The
relation $\eta_4=\eta_1\eta_3$ (from $\psi_4=\psi_1\psi_3$) makes the three
constraints for $j\in\{1,3,4\}$ multiplicatively dependent: any solution of the
$j=1$ and $j=3$ constraints automatically solves $j=4$. A common solution exists
iff the prescribed ratios are consistent powers of $\eta_j$, and when it does the
locus of $\tau\in\Z_3^2$ realizing it is a coset of $\ker\psi_3\cap\ker\psi_4$.
Since $\psi_3$ and $\psi_4$ separate the points of $\Z_3^2$ (their kernels
$\langle(1,2)\rangle$ and $\langle(1,1)\rangle$ meet only in $0$), the locus is a
single multiple of $(0,1)$.
\end{proof}

For quick reference we record the basepoint $C$-table on its own.

\begin{center}
\begin{tabular}{lcccc}
\toprule
 & $1$ & $\sx$ & $\sy$ & $\sx\sy$ \\
\midrule
$C_j([1])$ & $0$ & $\eta_j$ & $1$ & $\eta_j^2$ \\
\bottomrule
\end{tabular}
\end{center}

\noindent
Other basepoints $s_\star$ are obtained from this row by layer translation, and
the exponents are independent of the component $j$.

\subsection{The layer dictionary}

The cover argument also needs minimum weights on a single $\Z_3^2$-fiber, as a
function of how many nontrivial Frobenius orbits a Fourier support touches. We
record these in a small dictionary. Throughout, ``orbit set $W$'' means a set of
Frobenius orbits of characters of $\Z_3^2$, $n$ is the number of \emph{nontrivial}
orbits in $W$, and $\aug$ records whether the trivial orbit lies in $W$.

\begin{remark}[support-containment convention]
$d_3(W)$ denotes the minimum Hamming weight over all nonzero $f\in\F2[\Z_3^2]$
whose Fourier support is \emph{contained in} $W$ --- not equal to $W$. The two
notions genuinely differ: for instance the exact-support minimum at
$(n,\aug)=(2,T)$ is $5$, whereas $d_3(2,T)=3$, attained by a line whose Fourier
support is a proper subset of $W$. Every application below uses only the safe,
monotone form ``$\supp_{\mathrm{Fourier}}f\subseteq W\implies\wt f\ge d_3(W)$''.
\end{remark}

\begin{lemma}[layer dictionary]\label{lem:d3}
For a nonzero $f\in\F2[\Z_3^2]$ with Fourier support contained in an orbit set
$W$, the minimum weight $d_3(W)$ depends only on the pair
$(n,\aug)=(\#\,\text{nontrivial orbits in }W,\ \text{trivial}\in W)$, and is given
by the following table.

\begin{center}
\begin{tabular}{ccc}
\toprule
$n$ & $\aug$ & $d_3(W)$ \\
\midrule
$0$ & $T$ & $9$ \\
$1$ & $F$ & $6$ \\
$1$ & $T$ & $3$ \\
$2$ & $F$ & $4$ \\
$2$ & $T$ & $3$ \\
$3$ & $\cdot$ & $2$ \\
$4$ & $F$ & $2$ \\
$4$ & $T$ & $1$ \\
\bottomrule
\end{tabular}
\end{center}

In particular the only facts used downstream are: a weight-$1$ layer is a
$\delta$-point with full Fourier support, taking the values $\psi_j(t)$ across
the components; $d_3(\{1\})=d_3(\{3\})=6$; and $d_3(\{1,3\})=4$.
\end{lemma}

\begin{proof}
The group $\Z_3^2$ has $9$ elements and $9$ characters, organized into the
trivial orbit and four nontrivial Frobenius orbits $\{1,2,3,4\}$ (each of size
$2$, since Frobenius pairs $\psi$ with $\psi^2$). For a nonzero $f$ the Fourier
support is a nonempty union of orbits contained in $W$, and the constraint
``Fourier support $\subseteq W$'' is exactly the linear condition that the Fourier
coefficients off $W$ vanish, i.e.\ $f$ lies in a coset-free subspace whose
dimension is the number of characters in $W$.

We treat the rows by direct inspection of these small subspaces; each row is a
minimum over an explicit $\F2$-subspace of $\F2[\Z_3^2]$, all of dimension at most
$9$, and can be surveyed by hand.

\emph{Row $(0,T)$.} The Fourier support is contained in the trivial orbit alone,
so $f$ is forced to be constant on $\Z_3^2$; the only nonzero such $f$ is the
all-ones function, of weight $9$.

\emph{Rows with the trivial orbit absent ($\aug=F$).} Here $f$ has zero
augmentation, so $\wt f$ is even and $\ge 2$. For $n=1$ the admissible $f$ lie in
the $2$-dimensional space spanned by one nontrivial orbit; over $\F2$ this is the
set of indicator functions of the two cosets of $\ker\psi$ other than $\ker\psi$
itself together with their sum, and the minimum nonzero weight is $6$ (a single
nontrivial coset, of size $3$, does not have its Fourier support inside one orbit;
the smallest admissible configuration is a union of two cosets, weight $6$). For
$n=2$ two orbits give a $4$-dimensional space; the lightest admissible nonzero
function is a ``line'' (a coset of a $\ker\psi$), of weight $3$, but lines have a
nonzero trivial Fourier coefficient, so the lightest \emph{augmentation-zero}
admissible function has weight $4$. For $n=3$ the smallest weight drops to $2$ (a
difference of two points whose Fourier support omits exactly one orbit), and for
$n=4$ likewise $2$.

\emph{Rows with the trivial orbit present ($\aug=T$).} Adding the trivial orbit
adds the constant functions, allowing odd weights. A single nontrivial orbit plus
the trivial orbit ($n=1,\ \aug=T$) admits a coset of $\ker\psi$, weight $3$;
similarly $(2,T)$ admits a line of weight $3$. With all four nontrivial orbits and
the trivial orbit ($4,T$) the full algebra $\F2[\Z_3^2]$ is available and a single
point ($\delta$-function) of weight $1$ is admissible, with full Fourier support.

\emph{The downstream facts.} A weight-$1$ layer is a single point $\delta_t$; its
Fourier transform is $\psi_j(t)\ne0$ on every component, so its Fourier support is
all five orbits and its component values are exactly $\psi_j(t)$. Reading the
table at $(n,\aug)=(1,F)$ gives $d_3(\{1\})=d_3(\{3\})=6$ (a single nontrivial
orbit, no trivial orbit), and at $(2,F)$ gives $d_3(\{1,3\})=4$.
\end{proof}
\section{Theorem A: no small base cycles}\label{sec:thmA}

By the duality Lemma~\ref{lem:duality} it suffices to argue on the $Z$ side. We prove
Theorem~A and then exhibit a weight-$6$ logical, giving $d(\mathrm{base})=6$.

\begin{proof}[Proof of Theorem~A]
Suppose $u=(u_L,u_R)$ is a nonzero $1$-cycle of the base code with
\[
  A\cdot u_L=B\cdot u_R=:\sigma,\qquad \wt{u}=\wt{u_L}+\wt{u_R}\le 5 .
\]
By the parity congruence~\eqref{eq:PAR} we have $\wt{u_L}\equiv\wt{u_R}\pmod 2$, since
$\wt{\sigma}=\wt{A\cdot u_L}\equiv\wt{u_L}$ and $\wt{\sigma}=\wt{B\cdot u_R}\equiv\wt{u_R}$.
This eliminates every split of odd total parity mismatch, namely
\[
  (1,2),\ (2,1),\ (2,3),\ (3,2),\ (1,4),\ (4,1).
\]
The surviving splits with $\wt{u}\le 5$ are the one-sided splits $(k,0)$ and $(0,k)$, and
the two-sided splits $(1,1)$, $(1,3)$, $(3,1)$, $(2,2)$. We dispatch each.

\medskip
\noindent\textbf{One-sided splits $(k,0)$ and $(0,k)$.}
Here $u_R=0$ forces $A\cdot u_L=0$, i.e.\ $0\ne u_L\in\Ann(A)$ (and symmetrically
$0\ne u_R\in\Ann(B)$ for $(0,k)$). The Engine Lemma~\ref{lem:engine} gives the minimum
weight over the nonzero annihilators directly: every nonzero element of $\Ann(A)$ and of
$\Ann(B)$ has weight $\ge 6$. Indeed, on the $\F4$ layer frame the unit component of an
annihilator vanishes ($\hat z_0=\hat z_2=0$), so each surviving layer carries even weight;
a nonzero annihilator must activate a radical component, which has co-point-or-full support
and hence occupies at least three layers. Three nonzero layers of even weight force weight
$\ge 6$. This kills all one-sided splits with $k\le 5$.

\medskip
\noindent\textbf{Split $(1,1)$.}
Write $u_L=g$ and $u_R=r$ as single points. Then
$A\cdot g=x^a y^b\cdot A$ is the translate of $\supp A$ by $g$, a $3$-element set, and
likewise $B\cdot r$ is the translate of $\supp B$ by $r$. The cycle condition
$A\cdot g=B\cdot r$ forces these two translated $3$-sets to coincide. Two coinciding
$3$-sets have equal difference sets, so $\mathrm dA=\mathrm dB$. But
\[
  \mathrm dA=\{(0,\pm1),(3,\pm1),(3,\pm2)\},\qquad
  \mathrm dB=\{(\pm1,0),(\pm1,3),(\pm2,3)\}
\]
are disjoint and nonempty, a contradiction.

\medskip
\noindent\textbf{Splits $(1,3)$ and $(3,1)$.}
Consider $(1,3)$: $u_L=g$ a single point and $z:=\supp u_R$ a $3$-set. Then
$\wt{A\cdot g}=3$ exactly, so the cycle condition requires $\wt{B\cdot z}=3$. We compute
$\wt{B\cdot z}$ by inclusion--exclusion over the three translates $x^c y^d\cdot B$,
$x^c y^d\in z$. Each translate has weight $3$; by the bound $\mathrm{ov}\le1$ any two
distinct translates of $\supp B$ overlap in at most one cell, and a cell common to all
three is counted once. Writing $p$ for the number of overlapping (ordered-into-unordered)
column pairs ($0\le p\le 3$) and $T\in\{0,1\}$ for the number of cells common to all three
translates,
\[
  \wt{B\cdot z}=9-2p+4T .
\]
The only solution of $9-2p+4T=3$ with $p\le3$, $T\le1$ is $(p,T)=(3,0)$: all three pairs
overlap and there is no common triple cell. Thus $z=\{z_0,\ z_0+a,\ z_0+b\}$ is a
\emph{$\mathrm dB$-triangle}, meaning $a,\ b,\ b-a\in\mathrm dB$ and the three pairwise
overlap cells are distinct.

We enumerate. Up to translation we may take $z_0=0$; we need an unordered pair
$\{a,b\}\subset\mathrm dB$ with $b-a\in\mathrm dB$. Running over the six elements of
$\mathrm dB$ and the dozen candidate ordered pairs, the closure condition $b-a\in\mathrm dB$
selects exactly one triangle up to translation and reflection. Its two chiralities are
\[
  T_{+}=\{0,\ (1,0),\ (2,3)\},\qquad T_{-}=\{0,\ (1,0),\ (5,3)\}.
\]
For $T_{+}$ the three overlap cells coincide, giving $\wt{B\cdot z}=7\ne 3$, so $T_{+}$ is
excluded. For $T_{-}$ one computes $\wt{B\cdot z}=3$, with $B\cdot z$ a translate of
$y^3(1+x^2+x^4)$ --- a set of \emph{constant} $y$-coordinate. But the left side
$\sigma=A\cdot g=x^a y^b\cdot A$ has $y$-coordinates $b+\{0,1,2\}$ (the $y$-parts of
$\supp A=\{x^3,\ y,\ y^2\}$), which are pairwise distinct. So $\sigma$ cannot have constant
$y$-coordinate, and $T_{-}$ is excluded too.

The mirror split $(3,1)$ is identical with the roles of $A,B$ interchanged: $\supp u_L$ is a
$\mathrm dA$-triangle whose image $A\cdot u_L$ has constant $x$-coordinate, while
$\sigma=B\cdot r$ has the three pairwise distinct $x$-coordinates of $\supp B$. Contradiction.

\medskip
\noindent\textbf{Split $(2,2)$.}
Write $u_L=\{\ell_1,\ell_2\}$ and $u_R=\{r_1,r_2\}$. Introduce the two coordinate
projections
\[
  \pi_x,\pi_y\colon\F2[\Z_6^2]\to\F2[\Z_6],
\]
where $\pi_y$ sums over the $x$-coordinate (leaving a polynomial in $y$) and $\pi_x$ sums
over the $y$-coordinate. Each is a ring homomorphism, and
\[
  \pi_y(A)=1+y+y^2,\quad \pi_y(B)=y^3,\qquad
  \pi_x(A)=x^3,\quad \pi_x(B)=1+x+x^2 .
\]
Since $\mathrm{ov}\le1$, the weight of a sum of two distinct translates of $\supp A$ (resp.\
$\supp B$) is $6$ if the translates are disjoint and $4$ if they overlap. Hence
$\wt{\sigma}\in\{4,6\}$ computed from either side, and the two sides must agree.

\emph{Case $\wt{\sigma}=4$ (both pairs overlap).} Then $\ell_1-\ell_2\in\mathrm dA$ and
$r_1-r_2\in\mathrm dB$. Set $\delta_L=\delta_{\ell_1}+\delta_{\ell_2}$ and
$\delta_R=\delta_{r_1}+\delta_{r_2}$. We match $\pi_y$. On the left,
$\pi_y(\sigma)=(1+y+y^2)\cdot\pi_y(\delta_L)$, where $\pi_y(\delta_L)$ has $y$-gap equal to
the $y$-component of $\ell_1-\ell_2\in\mathrm dA$, namely $1$ or $2$; multiplying
$1+y+y^2$ by a $y$-gap-$1$ pair gives weight $2$, and by a $y$-gap-$2$ pair gives weight $4$.
On the right, $\pi_y(\sigma)=y^3\cdot\pi_y(\delta_R)$, where $\pi_y(\delta_R)$ has $y$-gap
equal to the $y$-component of $r_1-r_2\in\mathrm dB$, namely $0$ or $3$; this has weight $0$
(gap $0$, the two points coincide after projection) or $2$ (gap $3$). The only common value
is weight $2$, forcing the $\ell$-difference to have $y$-gap $1$ (so it is $(0,\pm1)$ or
$(3,\pm1)$) and the $r$-difference to have $y$-gap $3$.
\begin{itemize}
\item If $\ell_1-\ell_2=(0,\pm1)$ then $\pi_x(\delta_L)=0$, so $\pi_x(\sigma)=0$, i.e.\
  $(1+x+x^2)\cdot\pi_x(\delta_R)=0$ with $\wt{\pi_x(\delta_R)}\le2$. The annihilator of
  $1+x+x^2$ in $\F2[\Z_6]$ is the ideal generated by $(1+x)(1+x^3)=1+x+x^3+x^4$, whose
  minimum nonzero weight is $4$. A weight-$\le2$ element of this ideal is therefore $0$, so
  $\pi_x(\delta_R)=0$; that means $r_1,r_2$ share an $x$-coordinate and the $r$-difference is
  $(0,3)$, which is not in $\mathrm dB$. Contradiction.
\item If $\ell_1-\ell_2=\pm(3,1)$, then $\pi_x(\sigma)=x^3\cdot\pi_x(\delta_L)$ has weight $2$
  (the $x$-gap is $3\ne0$). Matching weight $2$ on the right,
  $\pi_x(\sigma)=(1+x+x^2)\cdot\pi_x(\delta_R)$, forces the $r$-pair to have $x$-gap $\pm1$,
  hence $r$-difference $\pm(1,3)$. Up to translation each side is now a single
  configuration: the left is $\sigma=A\cdot(1+x^3 y)\cdot t$ for some monomial $t$, with
  cells $\{x^3,\ y^2,\ x^3y^2,\ x^3y^3\}$ (here $x^3$ occurs with multiplicity), giving the
  $x$-coordinate multiplicity multiset
  \[
    \{3,1\}\qquad(\text{three cells at }x\text{-coordinate }3,\ \text{one at }0);
  \]
  the right is $\sigma=B\cdot(1+x y^3)\cdot t'$, with cells
  $\{y^3,\ x^2,\ x^2y^3,\ x^3y^3\}$ and $x$-coordinate multiplicity multiset
  \[
    \{2,1,1\}\qquad(\text{two cells at }x\text{-coordinate }2,\ \text{one each at }0,3).
  \]
  These multisets are invariant under translation, and $\{3,1\}\ne\{2,1,1\}$, so the two
  presentations of $\sigma$ are incompatible. Contradiction.
\end{itemize}

\emph{Case $\wt{\sigma}=6$ (both pairs disjoint).} Then $\ell_1-\ell_2\notin\mathrm dA$ and
$r_1-r_2\notin\mathrm dB$. We split on the $y$-gap of the $\ell$-difference.
\begin{itemize}
\item $y$-gap $0$ (so the $\ell$ $x$-gap is $\ne0$): then $\pi_y(\delta_L)=0$, so
  $\pi_y(\sigma)=0$, and $\pi_y(\sigma)=y^3\cdot\pi_y(\delta_R)=0$ forces the $r$-pair to
  have $y$-gap $0$ as well, hence $r$ $x$-gap $\ne0$. Matching
  $\wt{\pi_x(\sigma)}=2$ on the right, $(1+x+x^2)\cdot\pi_x(\delta_R)$ has weight $2$ only when
  the $r$ $x$-gap is $\pm1$, i.e.\ $r$-difference $(\pm1,0)\in\mathrm dB$ --- contradicting
  the disjoint hypothesis.
\item $y$-gap $\pm1$: then $\ell_1-\ell_2=(e,\pm1)$ with $e\in\{1,2,4,5\}$ (not $0,3$, else
  it lies in $\mathrm dA$). Now $\pi_y(\sigma)=(1+y+y^2)\cdot(\text{$y$-gap }1)$ has weight
  $2$, so matching the right $\pi_y(\sigma)=y^3\cdot\pi_y(\delta_R)$ forces the $r$ $y$-gap to
  be $3$, hence $r$-difference $(f,3)$ with $f\in\{0,3\}$ (else in $\mathrm dB$). Then on the
  left $\pi_x(\sigma)=x^3\cdot(\text{$x$-gap }e)$ has weight $2$, while on the right
  $\pi_x(\sigma)=(1+x+x^2)\cdot(\text{$x$-gap }f)$ has weight $0$ (if $f=0$) or $6$ (if
  $f=3$). No value is $2$. Contradiction.
\item $y$-gap $\pm2$ or $3$: then $\pi_y(\sigma)=(1+y+y^2)\cdot(\text{$y$-gap }2\text{ or }3)$
  has weight $4$ or $6$ on the left, while the right $y^3\cdot\pi_y(\delta_R)$ has weight
  $\le2$. No match. Contradiction.
\end{itemize}

All splits with $\wt{u}\le5$ are eliminated. Hence every nonzero $1$-cycle of the base code
has weight $\ge6$; by the duality Lemma~\ref{lem:duality} the same holds for $X$-type
$1$-cycles. In particular $d(\mathrm{base})\ge6$, and since nonzero stabilizers are nonzero
cycles, $\mu_Z,\mu_X\ge6$.
\end{proof}

\subsection*{Corollary A$'$: $d(\mathrm{base})=6$}

\begin{corollary}\label{cor:Aprime}
$d(\mathrm{base})=6$. Concretely, the weight-$6$ element
\[
  z^\ast=1+y+y^2+y^5+x^3+x^3y^4=(1+y+y^2+y^5)+x^3(1+y^4)\ \in\ \Ann(A)
\]
furnishes a $Z$-logical $(z^\ast,0)$ of weight $6$, and $\Ann(A),\Ann(B)$ have minimum
nonzero weight exactly $6$.
\end{corollary}

\begin{proof}
First, $A\cdot z^\ast=0$. Expanding $(x^3+y+y^2)\cdot z^\ast$ term by term,
\[
\begin{aligned}
  x^3\cdot z^\ast &= x^3 + x^3y + x^3y^2 + x^3y^5 + 1 + y^4,\\
  y\cdot z^\ast   &= y + y^2 + y^3 + 1 + x^3y + x^3y^5,\\
  y^2\cdot z^\ast &= y^2 + y^3 + y^4 + y + x^3y^2 + x^3,
\end{aligned}
\]
and summing over $\F2$ each of the nine distinct monomials
$1,\ y,\ y^2,\ y^3,\ y^4,\ x^3,\ x^3y,\ x^3y^2,\ x^3y^5$ appears exactly twice and so cancels.
Hence $A\cdot z^\ast=0$, i.e.\ $z^\ast\in\Ann(A)$, and $u^\ast:=(z^\ast,0)$ is a $1$-cycle of
weight $6$.

It remains to see that $u^\ast$ is not a stabilizer. A $Z$-stabilizer is
$\partial_2 w=(B\cdot w,\ A\cdot w)$ for some $w$; for it to have zero right block we need
$A\cdot w=0$, i.e.\ $w\in\Ann(A)$, and matching the left block forces $B\cdot w=z^\ast\ne0$.
Thus $w\in\Ann(A)\smallsetminus\ker\partial_2$ (a $w\in\ker\partial_2$ would give
$B\cdot w=0\ne z^\ast$). The one-block Lemma~\ref{lem:oneblock} then gives
$\wt{B\cdot w}\ge16$, contradicting $\wt{B\cdot w}=\wt{z^\ast}=6$. Hence no such $w$ exists, so $u^\ast$
cannot be a stabilizer. Therefore $u^\ast$ is a nontrivial $Z$-logical of weight $6$, giving
$d(\mathrm{base})\le6$; combined with Theorem~A this yields $d(\mathrm{base})=6$.

Finally, $z^\ast\in\Ann(A)$ has weight $6$ and Theorem~A's one-sided analysis shows every
nonzero element of $\Ann(A)$ has weight $\ge6$; the same statement for $\Ann(B)$ follows by
the duality Lemma~\ref{lem:duality}. Hence both annihilators have minimum nonzero weight
exactly $6$.
\end{proof}
\section{The double cover and Theorem~B}\label{sec:thmB}

The gross code is built from the base code by a free $\Z_2$ \emph{covering} construction
in the $x$-direction. This section fixes the covering notation --- the deck map $\sigma$,
the projection chain map $p$, the diagonal lift $\tau$, the two \emph{sheets}, and the
\emph{dangerous}/\emph{safe} dichotomy --- and uses it to prove the first weight floor,
Theorem~B. All of this notation is reused in the later sections.

\subsection{The free \texorpdfstring{$\Z_2$}{Z2} double cover in the \texorpdfstring{$x$}{x}-direction}

The base group is $G_{\mathrm b}=\Z_6\times\Z_6$ and the gross group is
$G_{\mathrm g}=\Z_{12}\times\Z_6$, with the \emph{same} polynomials $A=x^3+y+y^2$ and
$B=y^3+x+x^2$ in both. The quotient map $\Z_{12}\to\Z_6$, $x\mapsto x\bmod 6$, exhibits
$G_{\mathrm g}$ as a free $\Z_2$ cover of $G_{\mathrm b}$ in the $x$-coordinate: each
base point $(x_0,y)\in G_{\mathrm b}$ has the two lifts $(x_0,y)$ and $(x_0+6,y)$ in
$G_{\mathrm g}$. We call these the two \emph{sheets}: sheet $0$ is the set of gross points
with $x$-coordinate in $\{0,\dots,5\}$ and sheet $1$ the set with $x$-coordinate in
$\{6,\dots,11\}$. The free deck transformation is
\[
  \sigma\colon G_{\mathrm g}\to G_{\mathrm g},\qquad \sigma(x,y)=(x+6,\,y),
\]
an involution that swaps the two sheets and commutes with multiplication by $A$ and $B$
(both have $x$-support inside a single sheet's worth of shifts, and $\sigma$ is induced by
a group translation, so it is an $\F2[G_{\mathrm g}]$-module automorphism). Choosing sheet
$0$ as a fundamental domain identifies $\F2[G_{\mathrm g}]$ with $\F2[G_{\mathrm b}]^2$:
a gross element is written as a pair $(f_0,f_1)$ of base elements, $f_s$ being its
restriction to sheet $s$. Under this identification a gross $1$-chain is a pair of base
$1$-chains
\[
  v=(v_0,v_1),\qquad v_0,v_1\in\F2[G_{\mathrm b}]\oplus\F2[G_{\mathrm b}],
\]
and Hamming weight is additive across sheets:
\begin{equation}\label{eq:cover-weight}
  \wt{v}=\wt{v_0}+\wt{v_1}.
\end{equation}

\paragraph{Block form of the cover boundary.}
Multiplication by $A$ or $B$ on $\F2[G_{\mathrm g}]$ either keeps a monomial on its sheet
or moves it across the seam between the sheets, according to whether the relevant
$x$-shift wraps the $\Z_{12}$ coordinate past the sheet boundary. Splitting each base
boundary operator $\partial$ into its non-seam-crossing part $\dncut$ and its
seam-crossing part $\dcut$, so that
\[
  \partial=\dncut+\dcut\qquad\text{(an entrywise split on $\partial$),}
\]
the gross boundary acts on the sheet-pair $(v_0,v_1)$ by the $2\times2$ block matrix
\begin{equation}\label{eq:cover-block}
  \partial^{\mathrm{cov}}
  =\begin{pmatrix}\dncut & \dcut\\[2pt] \dcut & \dncut\end{pmatrix},
\end{equation}
the off-diagonal $\dcut$ blocks recording precisely the monomials that the deck map
carries from one sheet to the other. (The same block form holds for $\partial_1$ and
$\partial_2$; we suppress the subscript when the statement is uniform.) Equation
\eqref{eq:cover-block} is verified directly: it is the entrywise statement that each
nonzero entry of $\partial$ lands on the same sheet as its source or on the opposite
sheet, and the two cases are exactly $\dncut$ and $\dcut$.

\paragraph{The projection chain map.}
Summing the two block rows of \eqref{eq:cover-block} collapses the deck action: since
$\dncut+\dcut=\partial$, adding the rows sends $(v_0,v_1)$ to
$\partial v_0+\partial v_1=\partial(v_0+v_1)$. This is exactly the effect of the
\emph{projection}
\[
  p\colon \F2[G_{\mathrm g}]\to\F2[G_{\mathrm b}],\qquad p(v)=v_0+v_1,
\]
the $\F2$-linear map that adds the two sheets (equivalently, push-forward along the
covering $G_{\mathrm g}\to G_{\mathrm b}$). Applied blockwise to a $1$-chain,
$p(v_L,v_R)=(p(v_L),p(v_R))$, the displayed row-sum computation reads
\begin{equation}\label{eq:p-chainmap}
  \partial_1^{\mathrm{base}}\bigl(p(v)\bigr)
  =\text{(sum of the two block rows of }\partial_1^{\mathrm{cov}}v),
\end{equation}
so $p$ is a chain map: it sends gross $1$-cycles to base $1$-cycles. In particular,
\[
  v\in\ker\partial_1^{\mathrm{cov}}\ \Longrightarrow\ p(v)\in\ker\partial_1^{\mathrm{base}} .
\]
Moreover, because $p$ adds two $\F2$ vectors supported on disjoint sheets, no cancellation
can create support, only destroy it:
\begin{equation}\label{eq:p-weightdrop}
  \wt{p(v)}=\wt{v_0+v_1}\le\wt{v_0}+\wt{v_1}=\wt{v}.
\end{equation}

\paragraph{The diagonal lift.}
Dually, the diagonal map
\[
  \tau\colon\F2[G_{\mathrm b}]\to\F2[G_{\mathrm g}],\qquad \tau(u)=(u,u),
\]
is a chain map as well: applying \eqref{eq:cover-block} to $(u,u)$ gives in each block
$\dncut u+\dcut u=\partial u$, so $\tau$ commutes with the boundary and carries base
$1$-cycles to gross $1$-cycles. It is the section of $p$ on the diagonal,
$p\circ\tau=0$ over $\F2$ in the sense that $p(\tau(u))=u+u=0$; equivalently, $\img\tau$
is the kernel of $p$.

\paragraph{The kernel of \texorpdfstring{$p$}{p} is diagonal.}
If $p(v)=0$ then $v_0+v_1=0$, i.e.\ $v_1=v_0$, so $v=(v_0,v_0)=\tau(v_0)$ is
\emph{diagonal}. For such a $v$ the two rows of the cover cycle equation
$\partial_1^{\mathrm{cov}}v=0$ both read $(\dncut+\dcut)v_0=\partial_1^{\mathrm{base}}v_0=0$.
Hence
\begin{equation}\label{eq:diag-cycle}
  p(v)=0\ \text{ and }\ \partial_1^{\mathrm{cov}}v=0
  \ \Longleftrightarrow\ v=(v_0,v_0)\ \text{ with }\ \partial_1^{\mathrm{base}}v_0=0,
\end{equation}
and then $\wt{v}=2\wt{v_0}$ by \eqref{eq:cover-weight}.

\paragraph{Dangerous and safe classes.}
The induced map on homology $p_\ast\colon H_1(\mathrm{gross})\to H_1(\mathrm{base})$ splits
the logical classes of the gross code into two sectors, which organize the rest of the
proof:
\begin{itemize}
  \item a class is \emph{dangerous} if $p_\ast[v]=0$ (its representatives project to base
        \emph{boundaries}); and
  \item a class is \emph{safe} if $p_\ast[v]\ne0$ (its representatives project to nonzero
        base \emph{homology}).
\end{itemize}
Theorem~B treats both sectors at the level of cycles; the sharp dangerous-sector floor is
Theorem~C.

\subsection{Proof of Theorem~B}

\begin{thmB}[the cover floor]
The gross code $[[144,12,12]]$, the free $\Z_2$ double cover of the base in the
$x$-direction, satisfies $d(\mathrm{gross})\ge 6$. More precisely, every nonzero gross
$Z$-type $1$-cycle $v$ --- logical or stabilizer --- has $\wt{v}\ge 6$.
\end{thmB}

\begin{proof}
Let $v=(v_0,v_1)\ne 0$ be any gross $Z$-type $1$-cycle,
$\partial_1^{\mathrm{cov}}v=0$. We distinguish the two sectors according to the value of
the projection $p(v)=v_0+v_1$.

\emph{Safe case: $p(v)\ne 0$.} By \eqref{eq:p-chainmap}, $p(v)$ is a base $1$-cycle, and
it is nonzero by hypothesis. Theorem~A (in its strong form: \emph{no} nonzero base
$1$-cycle, whether logical or a boundary, has weight $\le 5$) gives $\wt{p(v)}\ge 6$.
Combining with the weight bound \eqref{eq:p-weightdrop},
\[
  \wt{v}\ \ge\ \wt{p(v)}\ \ge\ 6 .
\]

\emph{Dangerous case: $p(v)=0$.} By \eqref{eq:diag-cycle}, $v$ is diagonal,
$v=(v_0,v_0)$ with $v_0$ a base $1$-cycle, and $v_0\ne 0$ because $v\ne 0$. Theorem~A
applied to the nonzero base cycle $v_0$ gives $\wt{v_0}\ge 6$, whence
\[
  \wt{v}\ =\ 2\,\wt{v_0}\ \ge\ 12 .
\]

In either case $\wt{v}\ge 6$. Since a nontrivial $Z$-logical operator of the gross code is
in particular a nonzero $Z$-type $1$-cycle, every such operator has weight $\ge 6$, so
$d_Z(\mathrm{gross})\ge 6$; and $d_X(\mathrm{gross})=d_Z(\mathrm{gross})$ by the Duality
Lemma (Lemma~\ref{lem:duality}). Hence $d(\mathrm{gross})\ge 6$.
\end{proof}

\begin{remark}
The proof is a clean instance of the projection bound
$d(\mathrm{gross})\ge\min\{\,d(\mathrm{base}),\ \mu_Z(\mathrm{base})\,\}$: the safe case is
floored by $d(\mathrm{base})$ through $p$, and the dangerous case is floored by twice the
minimum cycle weight $\mu_Z(\mathrm{base})$ through the diagonal $\tau$. Theorem~A supplies
both inputs --- $d(\mathrm{base})\ge 6$ and $\mu_Z(\mathrm{base})\ge 6$ --- with no descent
through a tower of covers. The dangerous-case estimate $\wt{v}=2\wt{v_0}$ already suggests
the sharper bound $\ge 2\,d(\mathrm{base})=12$ pursued in Theorem~C, but at the level of
\emph{cycles} it only sees $\mu_Z$, since $v_0$ may there be a base \emph{boundary}.
\end{remark}
\section{The dangerous sector is heavy: Theorem~C}\label{sec:thmC}

Theorem~B caps the cover floor at $6$ because the safe sector does. The deeper fact,
and the route to the full distance, is that the \emph{dangerous} classes --- those in
the kernel of the projection $p_\ast\colon H_1(\mathrm{gross})\to H_1(\mathrm{base})$
--- carry weight at least $12=2\,d(\mathrm{base})$. The proof combines a slice identity
that reduces a dangerous cover cycle to a base stabilizer $b$ plus an off-support
excess, a classification of all base stabilizers of weight $\le 11$, and a lower bound
on the excess for each surviving $b$.

\subsection{The slice identity}\label{subsec:slice}

Recall from the double-cover construction (used for Theorem~B) that a cover $1$-chain is
a pair of base chains $v=(v_0,v_1)$, and that for each seam position $j$ the cover
boundary takes the block form
\[
  \partial_2^{\mathrm{cov}}=
  \begin{pmatrix}\dncut & \dcut\\[2pt]\dcut & \dncut\end{pmatrix},
  \qquad \partial_2=\dncut+\dcut,
\]
where $\dncut$ (non-crossing) and $\dcut$ (seam-crossing) split the base boundary
\emph{entrywise}, so $\supp(\dcut f)\subseteq\supp(\partial_2 f)$. The diagonal
embedding $\tau(u)=(u,u)$ sends base $1$-cycles to cover $1$-cycles, and the dangerous
cover cycles are exactly $\tau(Z_1)+\img\partial_2^{\mathrm{cov}}$, where $Z_1=\ker
\partial_1^{\mathrm{base}}$.

\begin{lemma}[slice identity]\label{lem:slice}
Let $v$ be a dangerous cover $Z$-cycle, write $b:=p(v)=\partial_2\,p(w)\in\Stab_Z(\mathrm{base})$
for the base stabilizer it projects to, and let $z_b$ be a base preimage with
$\partial_2 z_b=b$. Then, for every seam position $j$, the two sheets satisfy
$v_1=v_0+b$, and
\begin{equation}\label{eq:slice}
  \wt{v}=\wt{b}+2\,\bigl|\,v_0\ \text{off}\ \supp b\,\bigr|\ \ge\ \wt{b}+2\,m(b),
\end{equation}
where the \emph{off-support minimum} is
\[
  m(b):=\min\bigl\{\,\bigl|\,(\dcut z_b+u')\ \text{off}\ \supp b\,\bigr|
  \ :\ u'\in Z_1,\ [u']\notin\img\cmap\,\bigr\},
\]
and $\cmap\colon\ker\partial_2\to H_1(\mathrm{base})$, $\cmap[\zeta]=[\dcut\zeta]$, is the
Smith connecting map; its image $\img\cmap\subseteq H_1(\mathrm{base})$ is independent of the
seam and equals $\ker\tau_\ast$, the base classes whose diagonal lift $\tau(\cdot)=(\cdot,\cdot)$
is a cover boundary. Consequently Theorem~C is exactly the statement
\begin{equation}\label{eq:M}
  \text{(M)}\qquad \wt{b}+2\,m(b)\ \ge\ 12\qquad\text{for every }b\in\Stab_Z(\mathrm{base}).
\end{equation}
\end{lemma}

\begin{proof}
Since $v$ is dangerous it lies in $\tau(Z_1)+\img\partial_2^{\mathrm{cov}}$, so we may
write $v=\tau(u)+\partial_2^{\mathrm{cov}}w$ for some base cycle $u\in Z_1$ and cover
$2$-chain $w$. The projection of $\partial_2^{\mathrm{cov}}$ is $\partial_2$, hence
$p(v)=\partial_2\,p(w)=:b$ is a base stabilizer; fix any preimage $z_b$ with
$\partial_2 z_b=b$. Expanding the block boundary on $w=(w_0,w_1)$ and using
$\dncut+\dcut=\partial_2$, the two sheets of $v$ are
\[
  v_0=u+\dcut z_b+\partial_2 w_0,\qquad v_1=v_0+b
\]
up to the relabeling $w_0\leftrightarrow w_1$ of the two seam sheets; the second equation
is $v_0+v_1=p(v)=b$. The class $[v_0]$ in $H_1(\mathrm{base})$ equals
$[u]+[\dcut z_b]=[u]+\cmap[z_b]$, and since $[v]=\tau_\ast[v_0]$ in $H_1(\mathrm{cover})$, the
cover cycle $v$ is a nontrivial logical iff $[v_0]\notin\ker\tau_\ast=\img\cmap$, i.e.\ iff
$[\,u+\dcut z_b\,]\notin\img\cmap$ (equivalently $[u']\notin\img\cmap$ for the displayed
representative $u':=u+\dcut z_b$).

For the weight, set $x:=v_0$. Over $\F2$ the elementary identity
$\wt{x}+\wt{x+b}=\wt{b}+2\,\bigl|\,x\ \text{off}\ \supp b\,\bigr|$ holds coordinatewise
(on $\supp b$ the two terms contribute $1$ in total per cell; off $\supp b$ each cell of
$x$ contributes $2$). Since $\wt{v}=\wt{v_0}+\wt{v_1}=\wt{x}+\wt{x+b}$, this gives the
equality in \eqref{eq:slice}. The displayed representative of $v_0$ is
$\dcut z_b+u'$ with $u'$ a base cycle whose class avoids $\img\cmap$, so minimizing the
off-support part over all admissible $v$ with the same $b$ yields the bound
$\bigl|\,v_0\ \text{off}\ \supp b\,\bigr|\ge m(b)$. Finally, ranging over all $b$ and
substituting into \eqref{eq:slice}, the cover distance on the dangerous sector is
$\min_b\bigl(\wt{b}+2\,m(b)\bigr)$, so the bound \eqref{eq:M} is equivalent to
Theorem~C.
\end{proof}

\subsection{Structure of light base stabilizers}\label{subsec:structure}

Throughout this subsection $b=\partial_2 z=(B\cdot z,\ A\cdot z)$ is a base stabilizer,
and each block is decomposed into its four layers $s\in\Z_2^2$ via the layer frame, with
component transforms $V_j$ as in Lemma~\ref{lem:engine}. We record three structural
facts.

\begin{lemma}[parity]\label{lem:parity}
$\wt{b}$ is even.
\end{lemma}

\begin{proof}
The unit components agree, $\hA_0=\hB_0=1+u+v$, because $A$ and $B$ have the same
$s$-part multiset $\{1,s_x,s_y\}$; thus the two blocks have identical layer parities. By
\eqref{eq:PAR}, $\wt{B\cdot z}\equiv\wt z$ and $\wt{A\cdot z}\equiv\wt z\pmod 2$, so
$\wt{b}=\wt{B\cdot z}+\wt{A\cdot z}\equiv 2\wt z\equiv 0\pmod 2$.
\end{proof}

\begin{lemma}[the sharp one-block lemma]\label{lem:oneblock}
If $z'\in\Ann(A)\smallsetminus\ker\partial_2$ then $\wt{B\cdot z'}\ge 16$. Symmetrically,
if $z'\in\Ann(B)\smallsetminus\ker\partial_2$ then $\wt{A\cdot z'}\ge 16$. The bound is
attained.
\end{lemma}

\begin{proof}
We treat the $\Ann(A)$ side; the other is the mirror under the block-swapping
automorphism $A(y,x)=B(x,y)$. Since $z'\in\Ann(A)$, the unit components vanish,
$\hat z'_0=\hat z'_2=0$, and on the radical components $\hat z'_j\in(\hA_j)$ for
$j\in\{1,3,4\}$. We compute the value vectors of $B\cdot z'$ layerwise. On component $4$,
using $\hB_4=\omega\hA_4$ from the layer table,
\[
  V_4^B=\hB_4\,\hat z'_4=\omega\,\hA_4\,\hat z'_4=\omega\,V_4^A=0,
\]
since $z'\in\Ann(A)$. On component $3$, $V_3^B=\hB_3\,\hat z'_3$; as $\hat z'_3\in(\hA_3)$
and $\hB_3=\omega u+v$ is radical, the product lands in the socle $\F4\cdot uv$, hence
$V_3^B$ is a \emph{constant} vector. On component $1$, $V_1^B=\hB_1\,\hat z'_1$ with
$\hB_1=1+u+v$ a unit, so $V_1^B$ ranges over the ideal $(\hA_1)$, i.e.\ it is full,
co-point, or zero by the engine. The parity component is dead ($\hat z'_0=0$), so every
layer of $B\cdot z'$ has even weight, and the Fourier support of every layer is contained
in $\{1,3\}$ (orbits $1$ and $3$). The layer dictionary $\dncut$-costs of
Lemma~\ref{lem:d3} for such supports are
\[
  \{1\}\to 6,\qquad \{3\}\to 6,\qquad \{1,3\}\to 4.
\]
We now bound $\wt{B\cdot z'}=\sum_s\wt{(B\cdot z')_s}$ by cases on $V_3^B$ and $V_1^B$.

\smallskip
\noindent\emph{Case $V_3^B\ne 0$ (constant).} A nonzero constant on component $3$ is
present at \emph{every} layer, so all four layers are alive. The subcases on $V_1^B$:
\begin{center}
\renewcommand{\arraystretch}{1.2}
\begin{tabular}{lll}
\toprule
$V_1^B$ & per-layer supports & weight bound\\
\midrule
full      & four layers at $\{1,3\}$, cost $4$ each      & $4\cdot 4=16$\\
co-point  & three layers at $\{1,3\}$, one at $\{3\}$     & $3\cdot 4+6=18$\\
zero      & four layers at $\{3\}$, cost $6$ each         & $4\cdot 6=24$\\
\bottomrule
\end{tabular}
\end{center}

\smallskip
\noindent\emph{Case $V_3^B=0$.} Then component $1$ must be nonzero (else
$B\cdot z'=0$, putting $z'\in\ker\partial_2$, excluded). A nonzero $V_1^B$ has full or
co-point support, hence at least three live layers, each with support $\{1\}$ at cost $6$:
\[
  \wt{B\cdot z'}\ \ge\ 3\cdot 6=18.
\]

\smallskip
The minimum over all cases is $16$, attained by the full--$V_1^B$, nonzero-constant-$V_3^B$
configuration.
\end{proof}

\begin{lemma}[floor]\label{lem:floor-structure}
If $b\ne 0$ and $\wt{b}\le 10$, then \emph{both} blocks of $b$ have at least three nonzero
layers.
\end{lemma}

\begin{proof}
Suppose one block, say $A\cdot z$, has at most two nonzero layers; we derive
$\wt{b}\ge 12$. If $A\cdot z=0$ then $z\in\Ann(A)$, and $z\notin\ker\partial_2$ (else
$b=0$), so by Lemma~\ref{lem:oneblock} the other block has $\wt{B\cdot z}\ge 16$,
contradicting $\wt b\le 10$. If $A\cdot z\ne 0$, the engine support analysis on the
radical components shows that a single live block component already forces three live
layers: a nonzero radical ideal element is full or co-point (Lemma~\ref{lem:engine}),
which occupies $\ge 3$ layers, while the parity component contributes only even-weight
layers. Hence $A\cdot z$ cannot have a nonzero radical component with only two live
layers; the only way to keep $A\cdot z$ to $\le 2$ layers is to have all radical
components vanish, i.e.\ $z\in\Ann(A)$ again, returning to the previous case. Thus both
blocks carry $\ge 3$ nonzero layers whenever $0<\wt b\le 10$.
\end{proof}

\begin{lemma}[direction forcing]\label{lem:dirforce}
Suppose $f\in\img(A\cdot)$ has a zero layer and a weight-$2$ layer of the form
$\{p,\ p+\delta\}$ (together with some $\delta$-point layer). Then each radical component
$V_j$ is a nonzero co-point vector, and
\[
  V_j[s_{\mathrm{pair}}]=\psi_j(p)\bigl(1+\psi_j(\delta)\bigr)\ne 0
  \qquad(j\in\{1,3,4\}),
\]
so $\delta\notin\ker\psi_1\cup\ker\psi_3\cup\ker\psi_4$. The only direction left is the
$\ty$-direction; on the $B$ side the mirror forces the $\tx$-direction.
\end{lemma}

\begin{proof}
A zero layer forces each radical $V_j$ to a co-point (a full vector is nonzero on every
layer). On the weight-$2$ layer the value is $V_j[s_{\mathrm{pair}}]=\psi_j(p)+\psi_j(p+\delta)
=\psi_j(p)(1+\psi_j(\delta))$, which is nonzero exactly when $\psi_j(\delta)\ne 1$, i.e.\
$\delta\notin\ker\psi_j$. Using the kernels of the layer frame,
$\ker\psi_1=\mathrm{span}(1,0)$, $\ker\psi_3=\mathrm{span}(1,2)$,
$\ker\psi_4=\mathrm{span}(1,1)$, the only nonzero $\delta\in\Z_3^2$ avoiding all three is
$\delta\in\mathrm{span}(0,1)$, the $\ty$-direction.
\end{proof}

\subsection{Classification of light base stabilizers}\label{subsec:classification}

\begin{proposition}[light stabilizers]\label{prop:lightstab}
Every base stabilizer $b\in\Stab_Z(\mathrm{base})$ with $0<\wt b\le 11$ is one of
\begin{itemize}
\item the $36$ \emph{hexagons} $\partial_2\delta_g$, of weight $6$; or
\item the $216$ \emph{D-pairs} $\partial_2(\delta_g+\delta_{gd})$ with $d\in
  \mathrm dA\cup\mathrm dB$, of weight $10$.
\end{itemize}
In particular there is no base stabilizer of weight $8$, and the minimum nonzero
stabilizer weight is $\mu_Z=6$.
\end{proposition}

\begin{proof}
By Lemma~\ref{lem:parity} we have $\wt b\le 10$. By Lemma~\ref{lem:floor-structure} both
blocks have $\ge 3$ nonzero layers. The block-swapping automorphism $A(y,x)=B(x,y)$
exchanges the two blocks, so without loss of generality the lighter block is the $A$-block
$f=A\cdot z$, of weight $3$, $4$ or $5$. With $\ge 3$ nonzero layers, the layer profile of
$f$ is one of
\[
  (1,1,1)\ \mid\ (1,1,1,1),\ (2,1,1)\ \mid\ (2,1,1,1),\ (2,2,1),\ (3,1,1).
\]
We resolve each shape. In every surviving case $f$ is pinned inside a \emph{single
$\ty$-fibre}, and an endgame transfers the conclusion from the block $f$ to the full
stabilizer $b$. The comp-$1$ transfer operator
\[
  T:=\hA_1\,\hB_1^{-1}=\hA_1(1+u+v)=u+\omega v+(1+\omega)uv
\]
(here $\hB_1=1+u+v$ is a self-inverse unit) carries the $B$-block component-$1$ value
vector to the $A$-block one, $V_1^A=T\cdot V_1^B$, and satisfies $T\cdot\vec 1=0$; it is
the bookkeeping device for the kills below.

\paragraph{Profile $(1,1,1)$.} Three $\delta$-point layers; their radical components $V_j$
are co-points whose cross-layer ratios are $\eta$-powers with $j$-independent exponents
(the C-table of Lemma~\ref{lem:engine}). The separation of $\psi_3,\psi_4$ pins the three
cells to the difference pattern of $A\cdot\delta_g$ for some $g$. Endgame:
$z-\delta_g\in\Ann(A)$. The $A$-block of this profile has weight $3$ and $\wt b\le10$, so
$\wt{B\cdot z}=\wt b-\wt{A\cdot z}\le 10-3=7$; with $\wt{B\cdot\delta_g}=3$,
\[
  \wt{B(z-\delta_g)}\le\wt{B\cdot z}+\wt{B\cdot\delta_g}\le 7+3=10<16,
\]
so by Lemma~\ref{lem:oneblock} $z-\delta_g\in\ker\partial_2$, i.e.\ $z\equiv\delta_g$
modulo $\ker\partial_2$. Hence $b=\partial_2\delta_g$ is a \textbf{hexagon}.

\paragraph{Profile $(1,1,1,1)$.} All four $V_j$ are full, hence constant, so the four
$\delta$-cells coincide: $f$ is a $\delta$-column at a single fibre $t^\ast$. Any
completion to a light $b$ would need the $B$-block all-odd of weight $\le 6$, i.e.\ profile
$(1,1,1,1)$ or $(3,1,1,1)$. A single comp-$2$ transfer kills both at once. The $A$-side
pins
\[
  \hat z_2=\hA_2^{-1}\,\psi_2(t^\ast)\,uv=\psi_2(t^\ast)\,uv,
\]
since the unit $\hA_2$ fixes the socle ($u_0\cdot uv=\aug(u_0)\,uv=uv$). Then on the
$B$-block component $2$,
\[
  V_2^B=\hB_2\cdot\psi_2(t^\ast)\,uv=0,
\]
because the radical $\hB_2$ kills the socle. But either putative $B$-shape has a
$\delta$-point layer, where $V_2^B$ takes the nonzero value $\psi_2(\text{$t$-cell})$ ---
contradiction. \textbf{No light $b$.}

\paragraph{Profile $(2,1,1)$ --- the D-pair lemma.} Direction forcing
(Lemma~\ref{lem:dirforce}) puts the pair in the $\ty$-direction. Using $1+\eta=\eta^2$, the
C-ratio system has exactly one solution per pair-layer position, giving the three
single-fibre patterns realized by $A\cdot(\delta_g+\delta_{gd})$ for $d\in\{y,x^3y^2,x^3y\}$;
these are $108$ elements forming $3$ translation classes, and a direct check confirms they
exhaust the patterns produced by the ratio system. Since $\mathrm dA\cap\mathrm dB=\emptyset$,
a $\mathrm dA$-pair has block weights $(4,6)$. Endgame: a completion with $\wt b\le 10$ has
$\wt{B\cdot z}\le 6$, so
\[
  \wt{B(z-\text{pair})}\le\wt{B\cdot z}+\wt{B\cdot\text{pair}}\le 6+6=12<16,
\]
and Lemma~\ref{lem:oneblock} forces $z\equiv\text{pair}$ modulo $\ker\partial_2$. Hence
$b$ is the \textbf{D-pair}, $\wt b=10$. (This is precisely where the sharp bound $16$ is
needed: a weaker bound $\ge 12$ leaves a gap at exactly $12$.)

\paragraph{Profile $(3,1,1)$.} The weight-$3$ layer is either a line or a triangle. A
\emph{line} kills at least two of the three radical components at that layer --- only the
orbit orthogonal to the line survives --- contradicting the co-point support required of a
radical component. A \emph{triangle} $\{p,\ p+g,\ p+h\}$ has
$\kappa_j:=1+\psi_j(g)+\psi_j(h)$. The map $j\mapsto(j\!\cdot\! g,\ j\!\cdot\! h)$ is a
bijection, and $\kappa_j=0$ iff $(j\!\cdot\! g,\ j\!\cdot\! h)=\pm(1,2)$, so $\kappa_j=0$
for exactly one orbit class $j$. If that class is radical, the support is killed. If it is
the $A$-unit component $2$, the ratio system is solvable only when $\kappa_4=\kappa_1\kappa_3$
(from $\psi_4=\psi_1\psi_3$ and the $j$-independent C-exponents). The six triangle shapes
with dead orbit $2$ (taking $g_x=1,\ h_x=2$) are tabulated below, and each violates
$\kappa_4=\kappa_1\kappa_3$:
\begin{center}
\renewcommand{\arraystretch}{1.2}
\begin{tabular}{ccccc}
\toprule
$(g_y,h_y)$ & $\kappa_1$ & $\kappa_3$ & $\kappa_1\kappa_3$ & $\kappa_4$\\
\midrule
$(0,1)$ & $\omega$   & $\omega$   & $\omega^2$ & $1$\\
$(0,2)$ & $\omega^2$ & $\omega$   & $\omega^2$ & $\omega$\\
$(1,0)$ & $\omega$   & $\omega^2$ & $\omega$   & $\omega^2$\\
$(1,1)$ & $\omega^2$ & $\omega^2$ & $\omega^2$ & $\omega$\\
$(2,0)$ & $\omega^2$ & $\omega^2$ & $\omega$   & $1$\\
$(2,2)$ & $\omega$   & $\omega^2$ & $\omega$   & $\omega^2$\\
\bottomrule
\end{tabular}
\end{center}
In each row $\kappa_1\kappa_3\ne\kappa_4$, so the ratio system is unsolvable. Hence
$\img(A\cdot)$ has \textbf{no} $(3,1,1)$ element (and, by the mirror, neither does
$\img(B\cdot)$).

\paragraph{Profile $(2,1,1,1)$.} All four layers are alive. If the pair direction avoids
all radical kernels (the $\ty$-direction), all three radical $V_j$ are full, hence
constant, so the $\delta$-cells coincide at a fibre $t^\ast$ and the pair is
$\{t^\ast+e,\ t^\ast+2e\}$: one class. If the direction lies in one radical kernel, that
$V_j$ is a co-point while the other two force the $\delta$-cells equal --- but a co-point
takes three \emph{distinct} values on its support, a contradiction. To kill the surviving
class, observe a completion needs $\wt{B\cdot z}=5$ with three odd layers, so the
$B$-block is $(3,1,1)$ --- impossible by the mirror of the previous paragraph --- or
$(2,1,1,1)$, which the mirror classification pins to $\delta$-cells at a common fibre
$t_0$ and a $\tx$-pair. Since $\psi_1$ kills $\tx$-pairs,
\[
  V_1^B=\psi_1(t_0)\bigl(\vec 1+\delta_{s_P}\bigr),\qquad
  V_1^A=T\cdot V_1^B=\psi_1(t_0)\cdot\mathrm{shift}_{s_P}(T),
\]
a co-point vanishing at $s_P$. But the $A$-side says $V_1^A$ is a nonzero constant.
Contradiction. \textbf{No light $b$.}

\paragraph{Profile $(2,2,1)$.} Direction forcing puts both pairs in the $\ty$-direction;
the C-ratio bookkeeping pins the three layers to $\{t\}$, $\{t,t+e\}$, $\{t,t+2e\}$ (three
classes, single fibre). A completion needs $\wt{B\cdot z}=5$ with exactly one odd layer at
$s_\delta$; the only such profile with $\ge 3$ layers is $\{1,2,2\}$, so the $B$-block is a
mirror-$(2,2,1)$ with $\tx$-pairs. Then $V_1^B=\psi_1(t')\,\delta_{s_\delta}$ and
$V_1^A=T\cdot V_1^B$ is a co-point vanishing at $s_\delta$ --- but the $A$-side co-point
vanishes at $s_4\ne s_\delta$ and is nonzero at $s_\delta$. Contradiction.
\textbf{No light $b$.}

\smallskip
The only surviving shapes are $(1,1,1)$ (hexagons, weight $6$) and the $\mathrm dA$-branch
of $(2,1,1)$ (D-pairs, weight $10$). Counting: there are $36$ hexagons $\partial_2\delta_g$
(one per cell $g$, the group having $36$ elements) and $216$ D-pairs
$\partial_2(\delta_g+\delta_{gd})$, one for each unordered pair $\{g,g+d\}$ with
$d\in\mathrm dA\cup\mathrm dB$; since $\mathrm dA\cup\mathrm dB$ has $12$ elements and each
pair arises from the two choices $(g,d)$ and $(g+d,-d)$, this is $36\cdot12/2=216$. No
weight-$8$ stabilizer survives, and $\mu_Z=6$.
\end{proof}

\subsection{The off-support minima}\label{subsec:rungs}

For each light $b$ we now bound the off-support minimum $m(b)$ of
Lemma~\ref{lem:slice}.

\begin{lemma}[hexagon rung]\label{lem:rung-hex}
For a hexagon $b=\partial_2\delta_g$, $\ m(b)\ge 3$.
\end{lemma}

\begin{proof}
The seam split is entrywise, so $\supp(\dcut\delta_g)\subseteq\supp(\partial_2\delta_g)=
h(g)$, the hexagon. Hence
\[
  m(b)=\min\bigl\{\,\bigl|\,u'\ \text{off}\ h(g)\,\bigr|\ :\ u'\in Z_1,\
  [u']\notin\img\cmap\,\bigr\}.
\]
Suppose $\bigl|\,u'\ \text{off}\ h(g)\,\bigr|\le 2$. Replacing $u'$ by $u'+b$ if needed
(a base boundary, so the class $[u']$ is unchanged) we may also assume
$\bigl|\,u'\cap h(g)\,\bigr|\le 3$, because $h(g)$ has six
cells and $u'+b$ flips exactly those covered by $b$. Then $\wt{u'}\le 3+2=5$, so $u'$ is a
nonzero base cycle of weight $\le 5$ --- impossible by Theorem~A --- unless $u'=0$. But
$u'=0$ gives $[u']=0\in\img\cmap$, contradicting the admissibility of $u'$. Hence
$\bigl|\,u'\ \text{off}\ h(g)\,\bigr|\ge 3$.
\end{proof}

\begin{lemma}[D-pair rung]\label{lem:rung-dpair}
For a D-pair $b=\partial_2(\delta_g+\delta_{gd})$, $\ m(b)\ge 1$.
\end{lemma}

\begin{proof}
Write $b=b_1+b_2$ as the sum of the two hexagons $b_1=\partial_2\delta_g$,
$b_2=\partial_2\delta_{gd}$, so $\supp b=h\cup h'$ and $\supp(\dcut z_b)\subseteq h\cup h'
=\supp b\cup\{q^\ast\}$ for the single overlap cell $q^\ast$, giving an $11$-qubit union.
Suppose $m(b)=0$: some admissible cycle $u'$ (with $[u']\notin\img\cmap$) is supported
inside this $11$-qubit union $U=\supp b_1\cup\supp b_2$. Consider the four coset
representatives $u'+\{0,\ b_1,\ b_2,\ b_1+b_2\}$. At a qubit $q$ with
$(b_1(q),b_2(q))\ne(0,0)$, exactly two of the four values
$u'(q),\,u'(q)+b_1(q),\,u'(q)+b_2(q),\,u'(q)+b_1(q)+b_2(q)$ are nonzero, independently of
$u'(q)$; at a qubit with $(b_1(q),b_2(q))=(0,0)$ all four equal $u'(q)$, and $u'(q)=0$ there
since $\supp u'\subseteq U$. Hence the four weights sum to
\[
  2\,\bigl|\supp b_1\cup\supp b_2\bigr|=2\cdot 11=22<24=4\cdot 6,
\]
so at least one representative has weight $\le 5$, hence is the zero cycle by Theorem~A.
Then $u'\in\mathrm{span}\{b_1,b_2\}\subseteq\img\partial_2$, so $[u']=0\in\img\cmap$,
contradicting admissibility. Hence $m(b)\ge 1$.
\end{proof}

\subsection{Assembly of Theorem~C}\label{subsec:assembly}

\begin{proof}[Proof of Theorem~C]
We verify the factor-two bound \eqref{eq:M}, $\wt b+2\,m(b)\ge 12$, on every value of $b$.

\begin{itemize}
\item $b=0$. Here a representative $u'$ with $[u']\notin\img\cmap$ is a nonzero base
  cycle, so $m(0)\ge 6$ by Theorem~A, and $\wt b+2m(b)\ge 0+12=12$.
\item $b$ a hexagon. $\wt b=6$ and $m(b)\ge 3$ (Lemma~\ref{lem:rung-hex}), giving
  $6+2\cdot 3=12$.
\item $b$ a D-pair. $\wt b=10$ and $m(b)\ge 1$ (Lemma~\ref{lem:rung-dpair}), giving
  $10+2\cdot 1=12$.
\item $\wt b\ge 12$. Then $\wt b+2m(b)\ge 12$ trivially.
\end{itemize}
By Proposition~\ref{prop:lightstab} no other $b$ with $\wt b\le 11$ exists, so the four
cases above are exhaustive. Combining with the slice identity (Lemma~\ref{lem:slice}),
every nontrivial dangerous $Z$-logical $v$ satisfies
\[
  \wt v\ge\wt b+2\,m(b)\ge 12.
\]
By the Duality Lemma~\ref{lem:duality} the same bound holds on the $X$ side.

Tightness: a diagonal representative $\tau(u)=(u,u)$ over a weight-$6$ base logical $u$
with $[u]\notin\img\cmap$ has weight $12$, and such a $u$ exists --- the weight-$6$ logical
$u^\ast$ of Corollary~A$'$ satisfies $[u^\ast]\notin\img\cmap$, as shown in the proof of
Theorem~D. Hence the bound $12$ is attained.
\end{proof}
\section{The safe-sector reduction}\label{sec:safesector}

\subsection*{Route to Theorem~D}

We work on the $Z$ side throughout (the $X$ side follows by the Duality
Lemma~\ref{lem:duality}). Recall that the gross code is the free $\Z_2$ double
cover of the base in the $x$-direction: a cover $1$-chain is a pair $v=(v_0,v_1)$
of base chains (the two sheet coordinates), the projection $p(v)=v_0+v_1$ is a
chain map with $\wt v\ge\wt{p(v)}$, and the deck transformation
$\sigma\colon x\mapsto x+6$ acts on cover chains. We must show that every
nontrivial cover $Z$-logical $v$ has $\wt v\ge 12$.

Fix such a $v$ and consider the induced base class $[p(v)]\in H_1(\mathrm{base})$.
This splits the analysis into two sectors.

\begin{itemize}
\item \emph{Dangerous sector:} $[p(v)]=0$ in $H_1(\mathrm{base})$, i.e.\ the
class of $v$ lies in $\ker p_\ast$. Theorem~C handles this case directly:
every such nontrivial logical has weight $\ge 12$.
\item \emph{Safe sector:} $[p(v)]\ne 0$. Here $\wt v\ge\wt{p(v)}\ge\wt{u}$ for
the lightest cycle $u$ representing the nonzero base class $[p(v)]$, so a weight
floor on the relevant base classes transfers directly to $v$.
\end{itemize}

The present section reduces the safe sector to a single weight floor on the base
code, the statement we call \emph{(M-im)}: \emph{every base $1$-cycle whose class
lies in a nonzero Smith class has weight $\ge 12$} (Proposition~\ref{thm:Mim}). Granting
(M-im), the safe sector is immediate: the homotopy theorem~(R) below shows
$[p(v)]$ lands in the image of the Smith connecting map $\cmap$, so $[p(v)]\ne 0$
forces $\wt{p(v)}\ge 12$, whence $\wt v\ge 12$. Together with Theorem~C this gives
$d(\mathrm{gross})\ge 12$; Corollary~A$'$ and the tightness construction supply a
weight-$12$ logical, proving Theorem~D.

This subsection establishes the four structural facts on which (M-im) rests: the
homotopy theorem~(R), which confines the safe sector to the Smith classes; the
seam-flux functional, of which only the easy half is needed downstream; the
transport-and-parity description of the Smith cosets; and the confined CRT frame,
which kills two of the five component sectors outright and pins the remaining
three to two affine relations (the $\rho$-links).

\subsection{The homotopy theorem (R)}\label{thm:R}

The connecting map of the cover is the Smith map
\[
  \cmap\colon H_2(\mathrm{base})=\ker\partial_2\longrightarrow H_1(\mathrm{base}),
  \qquad \cmap[\zeta]=[\dcut\,\zeta],
\]
where $\dcut$ is the seam-crossing part of the cover boundary (so $\dcut\zeta$ is a base
$1$-cycle); its image $\img\cmap\subseteq H_1(\mathrm{base})$ is cut-independent and equals
$\ker\tau_\ast$, where $\tau_\ast\colon H_1(\mathrm{base})\to H_1(\mathrm{cover})$ is induced
by the diagonal transfer $\tau(u)=(u,u)$.

\begin{thmR}
Every cover cycle $v$ satisfies $[p(v)]\in\img\cmap$; equivalently
$\sigma_\ast=\mathrm{id}$ on $H_1(\mathrm{gross})$.
\end{thmR}

\begin{proof}
Work over $\F2[\Z_{12}\times\Z_6]$. Squaring $B=y^3+x+x^2$ kills its
$y$-dependence: since $(y^3)^2=y^6=1$ and the cross terms vanish in characteristic $2$,
$B^2=y^6+x^2+x^4=1+x^2+x^4$, so that
\[
  (1+x^2)\,B^2=(1+x^2)(1+x^2+x^4)=1+x^6 .
\]
Let $v=(v_L,v_R)$ be a cover cycle, so $A\cdot v_L=B\cdot v_R$, and set
\[
  z:=(1+x^2)\cdot B\cdot v_L\in\F2[\Z_{12}\times\Z_6].
\]
Using $\partial_2 z=(B\cdot z,\,A\cdot z)$, commutativity, and $A\cdot v_L=B\cdot v_R$,
\[
  \partial_2 z
  =\bigl((1+x^2)B^2\,v_L,\ (1+x^2)B\!\cdot\!A\,v_L\bigr)
  =\bigl((1+x^6)\,v_L,\ (1+x^6)\,v_R\bigr)
  =v+\sigma v,
\]
since $\sigma$ is multiplication by $x^6$. Thus $(1+\sigma)v=\partial_2 z$ is a
boundary, so $(1+\sigma)$ is null-homotopic on cycles and $\sigma_\ast=\mathrm{id}$
on $H_1(\mathrm{gross})$. Consequently
$\tau_\ast\circ p_\ast=(1+\sigma)_\ast=0$, i.e.\
$\img p_\ast\subseteq\ker\tau_\ast=\img\cmap$, which is the asserted inclusion.
\end{proof}

Only this \emph{inclusion} is load-bearing: the safe-sector floor needs
$[p(v)]\in\img\cmap\smallsetminus 0$, and $[p(v)]\ne 0$ is precisely the definition
of the safe sector. Neither the value $k=12$ nor the reverse inclusion of~(R) is
used in the distance bound.

\subsection{No double wrap and the seam flux}\label{lem:nodoublewrap}

Write $\dncut$ for the non-crossing part of the cover boundary, so
$\partial_2=\dcut+\dncut$, and let $\partial_1^{\mathrm c},\partial_1^{\mathrm{nc}}$
be the analogous splitting of $\partial_1$ at a fixed seam (cut) $j$.

\begin{lemma}[no double wrap]\label{lem:nodoublewrap-statement}
For every cut $j$,
\[
  \partial_1^{\mathrm c}\,\dcut=0,\qquad
  \partial_1^{\mathrm{nc}}\,\dncut=0,\qquad
  \partial_1^{\mathrm{nc}}\,\dcut=\partial_1^{\mathrm c}\,\dncut .
\]
\end{lemma}

\begin{proof}
An entry of $\partial_1\partial_2$ at (check $c$, face $f$) is a sum over two-step
paths $f\to\text{qubit}\to c$, one summand per factorization $c\,f^{-1}=a\cdot b$
along each route (left block: a $B$-step then an $A$-step; right block: an
$A$-step then a $B$-step). The routes contribute in pairs because $AB=BA$, so the
total number of paths at the entry is even.

The $x$-advance of a path is $D=s_x(a)+s_x(b)$. From the difference-set data,
$s_x$ of a single $A$- or $B$-step is at most $3$ and at most $2$ respectively, so
$D\le 3+2=5<6$. Moreover $D\equiv (c-f)_x\pmod 6$, and the constraint $0\le D\le 5$
pins $D$ to the unique residue: \emph{$D$ is the same integer for every path at
the entry}. A monotone path of total advance $D<6$ starting at $x_f$ crosses the
seam at most once, and whether it crosses is determined by the pair $(x_f,D)$
alone. Hence all paths at the entry share one crossing count.

If that count is $0$, every path lies entirely in the non-crossing part and the
even set of paths cancels inside $\partial_1^{\mathrm{nc}}\dncut$, while
$\partial_1^{\mathrm c}\dcut$ receives nothing. If the count is $1$, each path
crosses in exactly one of its two steps, so neither
$\partial_1^{\mathrm c}\dcut$ nor $\partial_1^{\mathrm{nc}}\dncut$ receives the
entry, and the paths distribute (with even total) between
$\partial_1^{\mathrm{nc}}\dcut$ and $\partial_1^{\mathrm c}\dncut$, forcing those
two entries equal. This proves all three identities.
\end{proof}

\begin{corollary}[seam flux]\label{cor:flux}
For each $\xi\in\ker H_X^{\mathsf T}$ the map
$\ell_\xi(w)=\xi^{\mathsf T}\partial_1^{\mathrm c}\,w$ descends to a well-defined
functional on $H_1(\mathrm{base})$, and $\img\cmap\subseteq\ker\ell_\xi$.
\end{corollary}

\begin{proof}
Well-definedness asks $\xi^{\mathsf T}\partial_1^{\mathrm c}\partial_2=0$. Since
$\partial_2=\dcut+\dncut$,
\[
  \xi^{\mathsf T}\partial_1^{\mathrm c}\partial_2
  =\xi^{\mathsf T}\partial_1^{\mathrm c}\dcut+\xi^{\mathsf T}\partial_1^{\mathrm c}\dncut
  =0+\xi^{\mathsf T}\partial_1^{\mathrm{nc}}\dcut ,
\]
using Lemma~\ref{lem:nodoublewrap-statement} ($\partial_1^{\mathrm c}\dcut=0$ and
$\partial_1^{\mathrm c}\dncut=\partial_1^{\mathrm{nc}}\dcut$). Now $\xi\in\ker H_X^{\mathsf T}$
means $\xi^{\mathsf T}\partial_1=0$, i.e.\
$\xi^{\mathsf T}\partial_1^{\mathrm c}=\xi^{\mathsf T}\partial_1^{\mathrm{nc}}$, hence
$\xi^{\mathsf T}\partial_1^{\mathrm{nc}}\dcut=\xi^{\mathsf T}\partial_1^{\mathrm c}\dcut=0$.
So $\ell_\xi$ kills boundaries and descends to $H_1$. For the inclusion, any class
in $\img\cmap$ is represented by $\dcut\zeta$ with $\zeta\in\ker\partial_2$, and
$\ell_\xi(\dcut\zeta)=\xi^{\mathsf T}\partial_1^{\mathrm c}\dcut\,\zeta=0$ by the
same first identity. Thus $\img\cmap\subseteq\ker\ell_\xi$.
\end{proof}

Only this easy inclusion $\img\cmap\subseteq\ker\ell_\xi$ is used below; it feeds
the tightness construction.

\subsection{The Smith cosets: transport and parity}\label{sec:smithcosets}

Each nonzero class of $H_2(\mathrm{base})=\ker\partial_2$ is realized by a cycle
$\zeta$, and the safe sector is governed by the cosets
\[
  C(\zeta)={\dcut}_0\,\zeta+\img\partial_2,
\]
taken at the fixed seam $j=0$. We first reduce the bookkeeping from all nonzero
classes to five orbit representatives.

\paragraph{Transport.}
The translations $T_x,T_y$ act on chains, and at the chain level one has the exact
identities
\[
  {\dcut}_j\circ T_x=T_x\circ{\dcut}_{j-1},\qquad
  {\dcut}_j\circ T_y=T_y\circ{\dcut}_j ,
\]
for every seam $j$. Combined with the cut-independence of $\cmap$, these give
$\cmap[T\zeta]=T\cdot\cmap[\zeta]$ for any translation $T$, so the connecting map
is translation-covariant. Consequently a weight floor proved at the level of
orbit representatives transports to the entire orbit, and to all classes the orbit
generates.

\begin{lemma}[dimension of $\ker\partial_2$]\label{lem:kerdim}
$\dim_{\F2}\ker\partial_2=6$; equivalently $\ker\partial_2$ has $2^6=64$ elements, of which
$63$ are nonzero.
\end{lemma}

\begin{proof}
Since $\partial_2\zeta=(B\cdot\zeta,\,A\cdot\zeta)$, we have
$\ker\partial_2=\Ann(A)\cap\Ann(B)$. The CRT layer frame splits
$\F2[\Z_6^2]\cong\F2[\Z_2^2]\times\F4[\Z_2^2]^4$ as $\F2$-algebras, and multiplication by
$A$ (resp.\ $B$) acts on component $j$ as multiplication by $\hA_j$ (resp.\ $\hB_j$). Hence
\[
  \ker\partial_2=\bigoplus_{j=0}^{4}\bigl(\Ann(\hA_j)\cap\Ann(\hB_j)\bigr),
\]
and each summand is read off the multiplier table and the Engine Lemma~\ref{lem:engine}:
\begin{itemize}[nosep]
\item $j=0$: $\hA_0=\hB_0=1+u+v$ is a unit, so $\Ann(\hA_0)=0$.
\item $j=1$: $\hB_1=1+u+v$ is a unit, so $\Ann(\hB_1)=0$ and the intersection is $0$.
\item $j=2$: $\hA_2=1+u+v$ is a unit, so $\Ann(\hA_2)=0$ and the intersection is $0$.
\item $j=3$: $\hA_3=u+\omega v$ and $\hB_3=\omega u+v$ are $\F4$-independent radicals, with
  $\Ann(\hA_3)=\F4\hA_3+\F4\,uv$ and $\Ann(\hB_3)=\F4\hB_3+\F4\,uv$ (Engine Lemma). An element
  $\alpha\hA_3+\beta\,uv=\gamma\hB_3+\delta\,uv$ of the intersection forces $\alpha=\gamma=0$
  (independence of $\hA_3,\hB_3$) and $\beta=\delta$, leaving $\F4\,uv$, of $\F2$-dimension $2$.
\item $j=4$: $\hB_4=\omega\hA_4$, so $\Ann(\hA_4)=\Ann(\hB_4)$ and the intersection is
  $\Ann(\hA_4)=\F4\hA_4+\F4\,uv$, of $\F2$-dimension $4$.
\end{itemize}
Summing, $\dim_{\F2}\ker\partial_2=0+0+0+2+4=6$.
\end{proof}

\begin{proposition}[orbit structure]\label{prop:smithorbits}
The set $\ker\partial_2\smallsetminus 0$ consists of exactly five translation
orbits, with $(\text{size},\text{weight})$ equal to
\[
  (9,16),\quad (12,18),\quad (36,18),\quad (3,24),\quad (3,24).
\]
Their orbit sums account for all $63$ nonzero Smith classes; in particular the
five canonical representatives, transported by $T_x,T_y$, cover every nonzero
Smith class.
\end{proposition}

\begin{proof}
By Lemma~\ref{lem:kerdim}, $\ker\partial_2$ has $63$ nonzero elements. The translation group
$\Z_6^2$ (order $36$) acts on $\ker\partial_2$ preserving Hamming weight, so computing the
weight of each nonzero element and grouping the $63$ of them into translation orbits is a
finite direct enumeration. It produces exactly five orbits; each orbit size is $36$ divided by
the order of the translation stabilizer of a representative, giving
$9=36/4$, $12=36/3$, $36=36/1$, $3=36/12$, $3=36/12$, with the stated minimum weights. As
$9+12+36+3+3=63$, these orbits exhaust the nonzero elements. Finally, translation-covariance of
$\cmap$ (above) gives $\cmap[T\zeta]=T\cdot\cmap[\zeta]$, so transporting the five
representatives by $T_x,T_y$ covers every nonzero Smith class.
\end{proof}

\noindent The five orbits are written $\mathrm{wt}\text{-}16$,
$\mathrm{wt}\text{-}18\mathrm a$, $\mathrm{wt}\text{-}18\mathrm b$,
$\mathrm{wt}\text{-}24\mathrm a$, $\mathrm{wt}\text{-}24\mathrm b$.

\paragraph{Parity.}
Two parity facts cut the candidate sub-$12$ weights down to a single residue. For
a layer $f\in\F2[\Z_3^2]$, weight is congruent to the trivial Fourier coefficient,
$\wt f\equiv\hat f(\mathrm{triv})\pmod 2$. Every $\zeta\in\ker\partial_2$ has even
columns: the relation $A\cdot\zeta=0$ reads columnwise as
$c_{i+3}=(y+y^2)c_i$, and $\aug(y+y^2)=0$, so each column of $\zeta$ has even
augmentation, hence even weight; the rows mirror this through $B$. Combining with
the previous fact, every element of every Smith coset has even Hamming weight (and
even seam value-cost). Therefore the only weights below $12$ that any coset element
could have are $6$, $8$, $10$, and the floor below rules out $6$ and $8$, leaving
weight exactly $10$ as the single case the kill step must eliminate.

\subsection{The confined frame: off-vanishing and the \texorpdfstring{$\rho$}{rho}-links}\label{ssec:offvanish}

We now diagonalize the relevant convolution operators by the CRT layer frame over
$\F4$, which decomposes $\F2[\Z_6^2]$ into five components $j=0,\dots,4$; write
$\hA_j,\hB_j$ for the layer multipliers. For a generic coset element
\[
  w={\dcut}_0\,\zeta+\partial_2 t,\qquad t\in\F2[\Z_6^2],
\]
the per-component data is
\[
  V_j(w)=\mathrm{off}_j+\bigl(\hB_j\,\hat t_j,\ \hA_j\,\hat t_j\bigr),
  \qquad \mathrm{off}_j=\mathrm{comp}_j({\dcut}_0\,\zeta),
\]
with the five components independent. Here $V_j=(V_jL,V_jR)$ records the left/right
block data at component $j$.

\begin{lemma}[off-vanishing at components $0,2$]\label{lem:offvanish-statement}
$\mathrm{off}_0=\mathrm{off}_2=0$ identically. Consequently the component-$0$ data of
every coset element is the diagonal pair $(V_0,V_0)$ with $V_0=\hB_0\hat t_0$ ranging
over all $16$ elements of $\F2[\Z_2^2]$ (as $\hB_0$ is a unit), and the component-$2$
data is the free graph $V_2L=\rho_2\cdot V_2R$, $\rho_2:=\hB_2\hA_2^{-1}$.
\end{lemma}

\begin{proof}
At components $0$ and $2$ the $A$-relation multiplier $\mathrm{comp}(y+y^2)$ equals
$Y:=1+s_y$: at $j=0$ both component characters are trivial, and at $j=2$ the
character is trivial on $t_y$. The column-collapse $v_i$ of $\zeta$ at such a
component therefore satisfies
\[
  v_i=Y\,v_{i+3}=Y^2 v_i=0,
\]
because $Y^2=(1+s_y)^2=1+s_y^2=0$ in the component algebra ($s_y^2=1$). The
seam-marked sums that build $\mathrm{off}_0$ and $\mathrm{off}_2$ are sums of these
vanishing collapses, so $\mathrm{off}_0=\mathrm{off}_2=0$. With
$\mathrm{off}_0=0$, the component-$0$ data is $(\hB_0\hat t_0,\hA_0\hat t_0)$; both
$\hA_0,\hB_0$ are units and at $j=0$ they coincide, giving the diagonal pair
$(V_0,V_0)$ with $V_0$ free over all of $\F2[\Z_2^2]$. With $\mathrm{off}_2=0$, the
component-$2$ data is $(\hB_2\hat t_2,\hA_2\hat t_2)$, i.e.\ the graph
$V_2L=\rho_2\,V_2R$ through the origin, $\rho_2=\hB_2\hA_2^{-1}$.
\end{proof}

\begin{lemma}[component $1$ collapses]\label{lem:c1zero}
The component-$1$ data satisfies the single relation $V_1R=\rho_1\,V_1L$ with
$\rho_1:=\hA_1\hB_1^{-1}$; in particular the residual obstruction $c_1$ vanishes
identically.
\end{lemma}

\begin{proof}
Let $c_0,\dots,c_5$ be the columns of $\zeta$ and $\hat u_i\in\F4[s_y]$ their
component-$1$ transforms. The crossing bookkeeping at seam $0$ gives
\[
  \mathrm{off}_1L=\hat u_4+\hat u_5+s_x\hat u_5,\qquad
  \mathrm{off}_1R=\hat u_3+s_x\hat u_4+\hat u_5 .
\]
The cycle relations transform to
\[
  \hat u_{i+3}=\tau\,\hat u_i\quad(\tau=\omega^2+\omega s_y,\ \text{a unit}),
  \qquad \hat u_{i-1}+\hat u_{i-2}=s_y\,\hat u_i,
\]
the second giving $\hat u_0=\hat u_1+s_y\hat u_2$ and hence
\begin{equation}\label{eq:D1}
  Y\hat u_1=\omega^2 Y\hat u_2.\tag{D1}
\end{equation}
The residual obstruction is
$c_1:=\mathrm{off}_1R+\rho_1^{-1}\!\cdot\!\mathrm{off}_1L$ (the right-block data
after normalizing the left through $\rho_1$). Reducing the group coefficients with
$\hB_1 X=s_x X$ and $\hB_1 Y=s_x Y$ (and the corresponding action on $\hA_1$), the
claim $c_1=0$ collapses to
\[
  Y\bigl[(X+\omega)\hat u_1+(\omega+\omega^2 s_x)\hat u_2\bigr]=0 .
\]
Applying~(D1) replaces $Y\hat u_1$ by $\omega^2 Y\hat u_2$, and the bracket
collapses to
\[
  Y(\omega^2+\omega+1)\hat u_2=0,
\]
which is identically zero since $\omega^2+\omega+1=0$ in $\F4$. Hence $c_1=0$, and
the component-$1$ graph passes through the origin: $V_1R=\rho_1 V_1L$.
\end{proof}

By Lemma~\ref{lem:offvanish-statement}, $\mathrm{off}_2=0$, so the component-$2$
graph likewise passes through the origin and the analogous obstruction $c_2$
vanishes by a one-line argument: with no offset, $V_2L=\rho_2 V_2R$ holds with no
constant correction.

\paragraph{The two $\rho$-links and the confined sets.}
Assembling the three preceding facts, on every Smith coset the component-$1,2$ data
obeys the two affine relations
\[
  V_1R=\rho_1\cdot V_1L,\qquad V_2L=\rho_2\cdot V_2R,
  \qquad \rho_1=\hA_1\hB_1^{-1},\ \ \rho_2=\hB_2\hA_2^{-1}.
\]
Relaxing the free sides, the \emph{confined sets} (the values available to the
dependent sides) are exactly $\img\rho_1$ and $\img\rho_2$, each of size $16$.
Indeed each $\rho_i$ satisfies
\[
  \rho_i^2=\aug(\rho_i)^2\cdot 1=0
\]
(squaring in $\F4[\Z_2^2]$ is the Frobenius-linear map and $g^2=e$ for
$g\in\Z_2^2$, so $\rho_i^2$ depends only on $\aug(\rho_i)$, which is $0$), while
$\rho_i\notin\F4\cdot\textstyle\sum_G$; hence $\img\rho_i$ is a $16$-element
$\F4$-subspace, the affine line of the labeling. These two $\rho$-links, together
with the free component-$0$ datum $V_0$ and the spine data at components $3,4$, are
the input to the slot frame, where the weight floor (M-im) is proved.
\section{The slot frame}\label{sec:slotframe}

All five Smith orbits of the base $Z$-stabilizer space are now placed in a
single coordinate system. Throughout, $b=\partial_2 z=(B\cdot z,\ A\cdot z)$ is a
base $Z$-stabilizer, decomposed into its five CRT components $V_0,\dots,V_4$ via
the layer frame of the Engine Lemma. The components $0$--$2$ are pinned by the
$\rho$-links (off${}_0=$ off${}_2=0$, the diagonal $V_0$-datum, and
$V_1^R=\rho_1\,V_1^L$, $V_2^L=\rho_2\,V_2^R$); the components $3$ and $4$ carry
the genuine freedom. The slot frame organizes that freedom.

\subsection{Slots, labelings, lines}

\emph{Slots} are the four elements of the $2$-part group $\Z_2^2=\{e,x,y,xy\}$
(the layers of the Engine Lemma). Component data live in $R=\F4[\Z_2^2]$,
identified with functions $\{\text{slots}\}\to\F4$, the \emph{slot values}.
Writing $X=1+s_x$, $Y=1+s_y$, $XY=\sum_{g\in\Z_2^2}g$, the slot values of
$a\cdot 1+\alpha X+\beta Y+\delta XY$ over $(e,x,y,xy)$ are
\[
  \bigl(a+\alpha+\beta+\delta,\ \alpha+\delta,\ \beta+\delta,\ \delta\bigr).
\]
The \emph{kill vector}
\[
  \kappa(v):=(a+\alpha+\beta,\ \alpha,\ \beta,\ 0)
\]
records the slot function of $v$ modulo the constant shift $\delta$: as the free
shift $\delta'$ varies, the zero set of $v+\delta'XY$ is a fibre of $\kappa(v)$.
Define the labeling vectors
\[
  \slab:=\kappa(\hB)=(\omega^2,\omega,1,0),\qquad
  \slab':=\kappa(\hA_3)=(\omega^2,1,\omega,0)=\slab\circ(x\leftrightarrow y),
\]
where $\hB=\hB_2=\hB_3=\hB_4=\omega X+Y$ and $\hA_1=\hA_3=X+\omega Y$, and the
two parity vectors
\[
  \theta:=(1,0,1,0),\qquad \tilde\theta:=(1,1,0,0).
\]

For brevity in the orbit-by-orbit analysis below we also write $m:=\slab$ and $m':=\slab'$
for these two labeling vectors.

\begin{lemma}[labeling facts]\label{lem:labelfacts}
The maps $\slab,\slab'\colon\{\text{slots}\}\to\F4$ are bijections;
$\tilde{\slab}:=\slab+\omega^2$ is an additive isomorphism
$\Z_2^2\xrightarrow{\ \sim\ }(\F4,+)$. Entrywise Frobenius gives
\[
  \kappa(\hA_4)=\slab'^{\,2}=\omega^2\slab,\qquad \slab^2=\omega^2\slab',
\]
and the two $\rho$-multipliers satisfy
\[
  \kappa(\rho_2)=\slab,\qquad \kappa(\rho_1)=\slab'.
\]
The parity vectors are the trace characters
$\theta=1+\operatorname{Tr}(\omega^2\tilde{\slab})$ and
$\tilde\theta=1+\operatorname{Tr}(\tilde{\slab})$.
\end{lemma}

\begin{proof}
Each assertion is a direct entrywise evaluation. From
$\slab=(\omega^2,\omega,1,0)$ one reads off that all four values are distinct,
so $\slab$ is a bijection, and likewise $\slab'$. Adding the constant
$\omega^2$ sends $\slab$ to $(0,\omega^2+\omega,\omega^2+1,\omega^2)$; over the
identification $\F4=\{0,1,\omega,\omega^2\}$ this is the standard additive chart
$\Z_2^2\to\F4$, so $\tilde{\slab}$ is an additive isomorphism. Squaring in
$\F4$ is the Frobenius automorphism $t\mapsto t^2$, which fixes $0,1$ and swaps
$\omega,\omega^2$; applying it entrywise to $\slab'=(\omega^2,1,\omega,0)$ gives
$(\omega,1,\omega^2,0)=\omega^2\slab$, and to $\slab$ gives
$(\omega,\omega^2,1,0)=\omega^2\slab'$. The identities $\kappa(\rho_2)=\slab$,
$\kappa(\rho_1)=\slab'$ are the kill vectors of $\rho_2=\hB_2\hA_2^{-1}$ and
$\rho_1=\hB_1\hA_1^{-1}$, computed from their unit/radical factors. Finally
$\operatorname{Tr}(\omega^2\tilde{\slab})$ and $\operatorname{Tr}(\tilde{\slab})$
are the $\F2$-valued slot functions $1+\theta$ and $1+\tilde\theta$, again by a
slotwise table.
\end{proof}

\paragraph{Confined and component lines.}
A direct check of the slot images of the confined images and the two
nondegenerate components yields the following descriptions. The free side of
each $\rho$-link contributes an unconstrained constant $c\in\F4$ (a free
$XY$-shift), so each set is a full affine line in $R$.

\begin{itemize}
\item \textbf{Confined lines.} The slot values of $\img\rho_2$ are
  $\{p\,\slab+c:p,c\in\F4\}$ and those of $\img\rho_1$ are
  $\{p\,\slab'+c\}$ --- the full affine line through the labeling direction.

\item \textbf{Component $3$.} On the left block,
  $V_3^L=\mathrm{off}_3^L+a_3\hB+\beta\,XY$ has slot values
  $\kappa_3^L+a_3\slab+c_3$ with $c_3\in\F4$ free (the coefficient $\beta$
  absorbs the $XY$-part of the offset) and $\kappa_3^L:=\kappa(\mathrm{off}_3^L)$.
  The right block is the mirror, with $\slab'$ in place of $\slab$ and
  $\kappa_3^R$. The spine coordinate $a_3$ is \emph{shared} between the two
  blocks; the two constants $c_3$ are free and independent.

\item \textbf{Component $4$.} The component-$4$ ideal is one-dimensional
  because $\hB_4=\omega\hA_4$, which forces the \emph{tie}
  \[
    V_4^L=\omega\,V_4^R+w_4,\qquad w_4:=\mathrm{off}_4^L+\omega\,\mathrm{off}_4^R,
  \]
  a fixed vector per orbit. In slot values the directions are
  \[
    k_4^L=\kappa_4^L+a_4\slab,\qquad k_4^R=\kappa_4^R+\omega^2 a_4\slab
  \]
  (the $\omega$-twist absorbs into $\omega a_4\slab'^{\,2}=a_4\slab$), with
  constants tied through a shared scalar $\gamma$:
  \[
    d_4^L=\omega\gamma+e_L,\qquad d_4^R=\gamma+e_R,
  \]
  where $e_\bullet$ denotes the $XY$-coefficient of the block's
  component-$4$ offset, and $e_L=\omega e_R+w_4^{XY}$ with $w_4^{XY}$ the
  $XY$-coefficient of $w_4$.
\end{itemize}

\paragraph{Per-orbit data.}
Computing the offsets $\mathrm{off}_3^\bullet,\mathrm{off}_4^\bullet$ of the
five canonical Smith cosets and reading off their kill vectors gives the table
below (offsets recorded as kill vectors; $\F4$ written $1,\omega,\omega^2$).

\begin{center}
\setlength{\tabcolsep}{5pt}
\begin{tabular}{lcccccc}
\toprule
orbit & $\kappa_3^L$ & $\kappa_3^R$ & $\kappa_4^L$ & $\kappa_4^R$ & $e_L$ & $e_R$\\
\midrule
$\text{wt-}16$  & $\omega\theta$ & $\theta$ & $\omega\theta$ & $\theta$ & $\omega$ & $\omega$\\
$\text{wt-}18a$ & $\slab$ & $\slab'$ & $(0,\omega,\omega^2,0)$ & $(\omega^2,\omega^2,\omega^2,0)$ & $\omega^2$ & $\omega^2$\\
$\text{wt-}18b$ & $(1,\omega,\omega^2,0)$ & $\omega^2\slab$ & $(\omega,\omega,1,0)$ & $(\omega,\omega^2,\omega,0)$ & $1$ & $1$\\
$\text{wt-}24a$ & $0$ & $0$ & $\omega\theta$ & $\theta$ & $\omega^2$ & $1$\\
$\text{wt-}24b$ & $\omega\theta$ & $\theta$ & $0$ & $0$ & $1$ & $\omega^2$\\
\bottomrule
\end{tabular}
\end{center}

\noindent The last row is the previous one with the component-$3$ and
component-$4$ offsets exchanged.

\subsection{The slot-cost rules}\label{ssec:slotcost}

We now record the per-slot layer cost. For a single slot $s$ and a single block,
fix the four component values $(v_0;v_{\mathrm{conf}},v_3,v_4)(s)$ at that slot
(here $v_{\mathrm{conf}}=v_2$ is the confined component); the block's other
unit-side component is free. The cost is the minimum layer weight realizing this
slot datum.

\begin{lemma}[slot-cost rules]\label{lem:slotcost}
The per-slot layer cost is governed by the radical rigidity
$\psi_2^2=\psi_3\psi_4$, $\psi_4=\psi_1\psi_3$ and the layer dictionary
$E\le 2$ value rigidity, and equals
\[
\begin{aligned}
  v_0=0:\quad & \text{$0$ alive}\to 0;\quad \text{$1$ alive}\to 4;\quad
    \text{$2$ alive}\to 2;\\
  & \text{$3$ alive}\to 2 \text{ if } T=1,\ \text{else } 4;\\
  v_0=1:\quad & \text{$3$ alive with $T=1$}\to 1\ (\text{$\delta$-point});\quad
    \text{else}\to 3,
\end{aligned}
\]
where ``alive'' counts the nonzero values among the block's three constrained
components and the rigidity datum is
\[
  T_L=v_2^2(v_3v_4)^{-1},\qquad T_R=v_4(v_1v_3)^{-1}.
\]
The $v_0$-free cost is the minimum over $v_0\in\{0,1\}$, namely
\[
  (0,\ 3,\ 2,\ 1\text{ if cheap else }3)\quad\text{by alive count }(0,1,2,3).
\]
\end{lemma}

\begin{proof}
A layer of $f$ supported at the slot $s$ contributes Fourier mass to the four
components according to the characters $\psi_1,\dots,\psi_4$; by the Engine
Lemma the relations $\psi_3=\psi_1\psi_2$ and $\psi_4=\psi_1\psi_3$ (hence
$\psi_2^2=\psi_3\psi_4$) constrain which value patterns are realizable, and the
layer dictionary fixes the minimum weight of a layer with $E\le 2$ prescribed
nonzero component values. The minimum layer weight realizing a slot datum with
$k$ alive constrained components and parity $v_0$ is then read off the dictionary
$d_3$: with $v_0=0$ (even layer), no alive components costs $0$; one alive
component is a single nonzero radical value, forcing a full or co-point layer of
weight $4$; two alive components costs $2$; three alive components costs $2$
precisely when the rigidity $T=1$ holds (so all three values are consistent with
a single $\delta$-line, attaining the dictionary minimum), and $4$ otherwise.
With $v_0=1$ (odd layer) the cheapest realization of three alive components with
$T=1$ is a single $\delta$-point, of weight $1$; every other odd configuration
costs $3$. Minimizing over the parity gives the stated $v_0$-free costs.
\end{proof}

\begin{remark}[slot parity]
Every per-slot cost is congruent to $v_0\pmod 2$, since an odd layer has odd
weight and an even layer even weight. Summing over the four slots, a block's
total cost is $\equiv\wt{V_0}\pmod 2$; hence any two blocks sharing the same
$V_0$ have costs of equal parity, and in particular every per-$(V_0,\gamma)$
cost sum is even.
\end{remark}

\subsection{Fibres of affine pencils}\label{ssec:pencil}

The slot-cost rules are driven by the alive set of a slot direction, i.e.\ by the
fibre partition of an affine slot function. This is governed by a single
elementary count.

\begin{lemma}[pair-ratio]\label{lem:pencil}
Let $k=\kappa+\lambda u$ with $u\colon\{\text{slots}\}\to\F4$ a bijection. For each
of the six unordered slot pairs $P=\{s,s'\}$ there is exactly one pencil parameter
\[
  \lambda_P=\bigl(\kappa(s)+\kappa(s')\bigr)\bigl(u(s)+u(s')\bigr)^{-1}
\]
at which $k(s)=k(s')$. The fibre partition of $k$ at a given $\lambda$ is read off
from $\{P:\lambda_P=\lambda\}$.
\end{lemma}

\begin{proof}
Because $u$ is a bijection, $u(s)+u(s')\ne 0$ for $s\ne s'$, so the equation
$\kappa(s)+\lambda u(s)=\kappa(s')+\lambda u(s')$ has the unique solution
$\lambda_P$. The slots $s,s'$ lie in the same fibre of $k$ exactly when
$\lambda=\lambda_P$, and collecting all pairs with $\lambda_P=\lambda$ gives the
fibre partition.
\end{proof}

\begin{corollary}[the component-$4$ trichotomy]\label{cor:comp4trich}
For the standard-form direction $k=b\slab+\omega\theta$ (so $\kappa=\omega\theta$,
$u=\slab$): $\lambda_P=0$ on the two $\theta$-constant pairs $\{e,y\}$,
$\{x,xy\}$; on the four $\theta$-split pairs
$\lambda_P=\omega\,\Delta\slab(P)^{-1}$, giving $\lambda=\omega$ on
$\{e,x\},\{y,xy\}$ and $\lambda=\omega^2$ on $\{e,xy\},\{x,y\}$. Consequently
$b\in\{0,\omega,\omega^2\}$ makes $k$ double-paired (its fibre partition is one of
the three pair-partitions of the slots, with alive sets of size $2$ or $4$),
while $b=1$ makes $k=(1,\omega,\omega^2,0)$ a bijection (alive sets of size $3$
with any dead slot).
\end{corollary}

\begin{proof}
Substitute $\kappa=\omega\theta=(\omega,0,\omega,0)$ and $u=\slab$ into the
pair-ratio formula. The pairs $\{e,y\}$ and $\{x,xy\}$ have equal
$\theta$-values, so $\kappa(s)+\kappa(s')=0$ and $\lambda_P=0$; the remaining
four pairs have $\kappa(s)+\kappa(s')=\omega$ and $\lambda_P=\omega(\Delta
\slab(P))^{-1}$, evaluated from the entries of $\slab$. The trichotomy is the
list of $b$-values for which $k$ is constant on a pair-partition versus
injective.
\end{proof}

\subsection{The chord-slope lemma and hyperbolic quadruples}\label{ssec:chord}

The cost of the most constrained configurations turns on whether four points of
the affine plane $AG(2,\F4)$ are in general position.

\begin{lemma}[chord slope]\label{lem:chord}
Let $u,k\colon\{\text{slots}\}\to\F4$ with $u$ bijective, fix a slot $z$, and set
\[
  g(s)=\bigl(u(s)+u(z)\bigr)\bigl(k(s)+k(z)\bigr)^{-1}
\]
on the three slots $s\ne z$ (defined when $k$ is injective off the fibre of $z$).
Then $g(s)=g(s')$ if and only if the three points $(u(s),k(s))$,
$(u(s'),k(s'))$, $(u(z),k(z))$ of $AG(2,\F4)$ are collinear. In particular, if
no three of the four points $(u(s),k(s))$ are collinear, then $g$ is injective
for every choice of $z$.
\end{lemma}

\begin{proof}
The quantity $g(s)$ is the reciprocal slope of the chord from $(u(z),k(z))$ to
$(u(s),k(s))$. Two chords from a common point $(u(z),k(z))$ have equal slope
exactly when the three endpoints are collinear. Thus $g(s)=g(s')$ iff the three
named points lie on a line, and $g$ is injective for all $z$ iff no three of the
four points are collinear.
\end{proof}

\begin{lemma}[hyperbolic quadruple]\label{lem:hyperbolic}
For $c\in\F4^\times$, the four points
\[
  H_c=\{(t,\ c\,t^{-1}):t\in\F4^\times\}\cup\{(0,0)\}
\]
have no three collinear, and every chord $(\Delta u,\Delta v)$ of $H_c$ satisfies
$\Delta u\cdot\Delta v=c$.
\end{lemma}

\begin{proof}
A line $v=\lambda u$ through the origin meets the hyperbola $uv=c$ where
$\lambda u^2=c$; since squaring is a bijection of $\F4$, this has exactly one
solution, so such a line meets $H_c$ in at most the origin and one further point.
A line $v=\lambda u+c'$ with $c'\ne 0$ meets the hyperbola where
$\lambda u^2+c'u+c=0$, a quadratic with at most two roots. Hence no line meets
$H_c$ in three points. For the chord identity, if $(u_1,c u_1^{-1})$ and
$(u_2,c u_2^{-1})$ are two finite points then $\Delta u=u_1+u_2$ and
$\Delta v=c(u_1^{-1}+u_2^{-1})$, so
\[
  \Delta u\cdot\Delta v=c(u_1+u_2)(u_1^{-1}+u_2^{-1})
   =c\bigl(1+u_1u_2^{-1}+u_2u_1^{-1}+1\bigr)
   =c\bigl(u_1u_2^{-1}+u_2u_1^{-1}\bigr),
\]
and over $\F4$ one checks $r+r^{-1}=1$ for every $r\in\F4^\times\smallsetminus\{1\}$
(here $r=u_1u_2^{-1}\ne 1$), giving $\Delta u\cdot\Delta v=c$. The chord from the
origin to $(u_1,cu_1^{-1})$ has $\Delta u\cdot\Delta v=u_1\cdot cu_1^{-1}=c$ as
well.
\end{proof}

\begin{remark}
The deepest slope case of the standard-form walk uses exactly one instance: the
pair $(\slab,\slab+\omega\theta)$ is the quadruple
$\{(\omega^2,1),(\omega,\omega),(1,\omega^2),(0,0)\}=H_{\omega^2}$, since the four
products $\slab\cdot(\slab+\omega\theta)$ equal $\omega^2,\omega^2,\omega^2,0$.
\end{remark}

\subsection{Cost-preserving moves and the standard form}\label{ssec:stdform}

\begin{lemma}[cost-preserving moves]
The block cost is invariant under each of the following operations.
\begin{enumerate}
\item \textbf{Slot relabelings} applied simultaneously to all components and to
  $V_0$ (the cost is a sum over slots).
\item \textbf{Translation scalings}
  $(v_{\mathrm{conf}},v_3,v_4)\mapsto(s_2 v_2,s_3 v_3,s_4 v_4)$ with
  $s_2^2=s_3 s_4$ (the nine cell symmetries of the value table); the confined
  line is scale-invariant.
\item \textbf{Frobenius} applied to all values (a value-table symmetry);
  moreover the right-block value table $M_2(v_0,\cdot,\cdot,v_4)$ equals
  $M_1(v_0,\cdot,\cdot,v_4^2)$, so an $R$-block is an $L$-type problem after
  Frobenius on its component-$4$ values.
\end{enumerate}
\end{lemma}

\begin{proof}
(1) The block cost is a sum over slots of per-slot costs, and a slot relabeling
permutes the summands. (2) The rigidity datum $T_L=v_2^2(v_3v_4)^{-1}$ is
unchanged by $(v_2,v_3,v_4)\mapsto(s_2 v_2,s_3 v_3,s_4 v_4)$ precisely when
$s_2^2=s_3 s_4$, and the confined line $\langle\slab\rangle$ is closed under
scaling; the alive sets and hence the per-slot costs are preserved. (3) Frobenius
is an automorphism of $\F4$ fixing the value-table relations, and the stated
identity $M_2=M_1(\cdot,\cdot,\cdot,v_4^2)$ is the entrywise comparison of the
left- and right-block rigidity data, which differ only by squaring the
component-$4$ value.
\end{proof}

\begin{definition}[standard form]
Let $S(a,b)$ be the $v_0$-free block problem with confined line
$\langle\slab\rangle$, component-$3$ direction $v_3=a\slab+c_3$, and
component-$4$ direction $v_4=b\slab+\omega\theta+c_4$, with $c_3,c_4\in\F4$ free.
\end{definition}

\begin{lemma}[reduction to standard form]\label{lem:stdform}
With the spine coordinate written $(a_3,a_4)$, the four $\text{wt-}24$ block
problems reduce as
\[
  L(24a)=S(a_3,a_4),\quad L(24b)=S(a_4,a_3),\quad
  R(24a)\cong S(a_3,a_4^2),\quad R(24b)\cong S(a_4^2,a_3).
\]
\end{lemma}

\begin{proof}
For $L(24a)$ this is the definition: the per-orbit table gives $\kappa_3^L=0$ and
$\kappa_4^L=\omega\theta$, so $v_3=a_3\slab+c_3$ and
$v_4=a_4\slab+\omega\theta+c_4$. For $L(24b)$ the offsets exchange roles, and the
$v_3\leftrightarrow v_4$ symmetry of $T_L=v_2^2(v_3v_4)^{-1}$ identifies the cost
with $S(a_4,a_3)$. For $R(24a)$: apply Frobenius to component $4$ (move 3), which
sends the data to $(\slab';\,a_3\slab';\,a_4^2\slab'+\theta)$ using $\theta^2=
\theta$ and $(\omega^2 a_4\slab)^2=a_4^2\slab'$; apply the slot swap
$x\leftrightarrow y$ (move 1) to obtain $(\slab;\,a_3\slab;\,a_4^2\slab+
\tilde\theta)$; apply the slot map $\sigma$ induced by
$\tilde{\slab}\mapsto\omega^2\tilde{\slab}$ (move 1), which sends
$\tilde\theta\mapsto\theta$ and $\slab\mapsto\omega^2\slab+1$; finally scale by
$(\omega,\omega,\omega)$ (move 2). The result is $S(a_3,a_4^2)$. The orbit
$R(24b)$ mirrors $R(24a)$ with the spine coordinates exchanged.
\end{proof}

\subsection{The achiever-structure lemma}\label{ssec:achiever}

Fix an orbit and a spine cell. For each shared pair $(V_0,\gamma)$ let
$\min_L(V_0,\gamma)$ and $\min_R(V_0,\gamma)$ be the per-block linked minima,
each minimized over the block's own knobs --- the confined point on its line, the
constant $c_3$, and the free side of its per-slot minimizations. By the slot
parity remark both minima are $\equiv\wt{V_0}\pmod 2$, so their sum is even. The
cell value is
\[
  m(\text{cell})=\min_{(V_0,\gamma)}\bigl(\min_L(V_0,\gamma)+\min_R(V_0,\gamma)\bigr).
\]

\begin{lemma}[achiever structure]\label{lem:achiever}
Suppose $m(\text{cell})\ge 10$ for every cell of the orbit. Then a weight-$10$
coset element must sit at some cell, in a configuration of cost exactly $10$ with
every slot at its minimum-cost value, with both free sides in the corresponding
per-slot argmin sets, and with both $\rho$-links satisfied. The set of such
cost-$10$ configurations is exactly
\[
  \bigcup_{\substack{(V_0,\gamma)\\ \min_L+\min_R=10}}
    \operatorname{Argmin}_L(V_0,\gamma)\times\operatorname{Argmin}_R(V_0,\gamma).
\]
\end{lemma}

\begin{proof}
The total weight of a coset element is the sum of its layer weights, which equals
its total block cost; since each slot cost is at least its minimum-cost value,
weight $10$ together with cost $\ge 10$ forces every layer to be a
minimum-weight layer for its slot datum (\emph{slot-exactness}: cost $=$ weight
means each layer attains its per-slot minimum). For the chosen cell, a cost-$10$
pair $(\mathrm{cost}_L,\mathrm{cost}_R)$ with $\mathrm{cost}_i\ge\min_i$ and
$\min_L+\min_R\ge 10$ (by hypothesis) forces $\mathrm{cost}_i=\min_i$ and
$\min_L+\min_R=10$. Hence the configuration realizes a $(V_0,\gamma)$ with
$\min_L+\min_R=10$, with each block in its argmin set, which is the displayed
union.
\end{proof}

\begin{remark}
The achiever-structure lemma reduces the weight-$10$ kill to two finite
sub-problems per cell: first, the \emph{cost-$8$ kill} $\min_L+\min_R\ge 10$ for
all $(V_0,\gamma)$ --- by the slot parity remark and the floors $\min\ge 3$, the
splits that could violate it are $(3,3)$, $(4,4)$, $(3,5)$ and $(5,3)$ --- and second, the explicit
\emph{loci} $\{(V_0,\gamma):\min_L+\min_R=10\}$ together with their argmin sets,
the \emph{achiever lists}. Each achiever is then killed against the
$\rho$-links.
\end{remark}
\section{The weight-\texorpdfstring{$24$}{24} orbits: the standard-form walk}\label{sec:wt24}

We now discharge the six weight-$24$ Smith classes of the dangerous sector. Recall the
slot frame: the four slots are the elements of $Z_2^2=\{e,x,y,xy\}$, component data are
functions slots~$\to\F4$, and the bijection $m=\kappa(\hB)=(\omega^2,\omega,1,0)$ together
with $\theta=(1,0,1,0)$ generate every line that occurs below. Each weight-$24$ block
contributes a configuration value $v_2$ on the \emph{configuration line} (the conf line
$\langle m\rangle$), together with two constrained component values $v_3,v_4$, and the
linked cost of the block is the sum over slots of the per-slot cost from the slot-cost
rule (Lemma~\ref{lem:slotcost}). The standard form of Lemma~\ref{lem:stdform} reduces all
four weight-$24$ block tables to a single two-parameter problem.

Recall (standard-form Definition, with $m=\slab$) that $S(a,b)$ is the \emph{$v_0$-free}
value of the standard-form block problem with conf line $\langle m\rangle$, comp-$3$ value
$v_3=am+c_3$, and comp-$4$ value $v_4=bm+\omega\theta+c_4$, minimized over the configuration
scale $p$, the constants $c_3,c_4\in\F4$, and the per-slot free unit-side components.
Explicitly, the per-slot cost is the $v_0$-free cost
\[
  \bigl(0,\ 3,\ 2,\ \text{$1$ if cheap else $3$}\bigr)
\]
indexed by the number of \emph{alive} (nonzero) values among the three constrained
components $(v_2,v_3,v_4)$ at that slot, ``cheap'' meaning
$T:=v_2^2(v_3v_4)^{-1}=1$ there.

\begin{proposition}\label{prop:wt24}
$S(a,b)\ge 6$ for all $(a,b)\in\F4^2$. Consequently every weight-$24$ block table is
$\ge 6$ everywhere, and since the $v_0$-free cost lower-bounds every cost with $V_0$
fixed, every weight-$24$ spine cell has linked value $\ge 6+6=12$. Thus the homological
floor (M-im) holds on the six weight-$24$ Smith classes.
\end{proposition}

This is the weight-$24$ part of Proposition~\ref{prop:ctablefloor}.

\begin{proof}
By Lemma~\ref{lem:stdform} the four weight-$24$ block tables are obtained from $S$ by
cost-preserving relabelings:
\[
  L(24a)\ \text{at}\ (a_3,a_4)=S(a_3,a_4),\quad
  L(24b)=S(a_4,a_3),\quad
  R(24a)\cong S(a_3,a_4^2),\quad
  R(24b)\cong S(a_4^2,a_3),
\]
so it suffices to prove $S(a,b)\ge 6$ for all $16$ pairs $(a,b)$. We organize the survey
by the support type of the configuration value $v_2$ and of $v_3$.

\medskip
\noindent\emph{Support trichotomy.}
Write the configuration value along its line as $v_2=pm+c_2$ with scale $p\in\F4$ and
constant $c_2\in\F4$. There are three support types:
\begin{itemize}[leftmargin=2em]
  \item \emph{dead}: $p=c_2=0$, so $v_2\equiv 0$;
  \item \emph{full constant}: $p=0\ne c_2$, so $v_2$ is a nonzero constant (alive on all
        four slots);
  \item \emph{co-point at $z_2$}: $p\ne 0$, so $v_2=p\,(m+m(z_2))$ vanishes at the unique
        slot $z_2$ with $m(z_2)=-c_2/p$ and is alive on the other three.
\end{itemize}
The comp-$3$ value $v_3=am+c_3$ is dead ($a=0=c_3$), full ($a=0\ne c_3$), or a co-point at
the unique slot $z_3$ ($a\ne 0$). The comp-$4$ value $v_4=bm+\omega\theta+c_4$ is governed
by the pencil trichotomy of Lemma~\ref{lem:pencil}: for $b\in\{0,\omega,\omega^2\}$ the
direction $bm+\omega\theta$ is \emph{double-paired}, so $v_4$ has an alive set $S_4$ of size
$|S_4|\in\{2,4\}$ ($v_4$ is constant on $S_4$ when $|S_4|=2$, and two-valued -- constant on
each fibre pair -- when $|S_4|=4$); for $b=1$ the direction is a bijection, so
$|S_4|=3$ with $v_4$ injective on $S_4$ and the dead slot $z_4$ free.

Throughout, a slot with $k$ alive components among $(v_2,v_3,v_4)$ costs $0,3,2,1/3$ for
$k=0,1,2,3$ respectively, the last being $1$ exactly when $T=v_2^2(v_3v_4)^{-1}=1$ at that
slot and $3$ otherwise (Lemma~\ref{lem:slotcost}). We split into $a=0$ (buckets
A1--A6) and $a\ne 0$ (buckets B1--B4), and in each bucket exhibit the minimizing
configuration explicitly.

\medskip
\noindent\textbf{Case $a=0$ (comp $3$ dead or full).}

\smallskip
\noindent\emph{A1 (comp $3$ dead, conf dead).}
Here $v_2\equiv 0$ and $v_3\equiv 0$, so every alive slot has only $v_4$ alive and costs
$3$; the $4-|S_4|$ dead slots cost $0$. The cost is $3|S_4|\ge 6$. The minima at
$|S_4|=2,3,4$ are $6,9,12$.

\smallskip
\noindent\emph{A2 (comp $3$ dead, conf co-point at $z_2$).}
Now $v_2$ is alive off $z_2$ and $v_3\equiv 0$. With $|S_4|=2$: if $z_2\notin S_4$ the two
slots of $S_4$ each have $v_2,v_4$ alive (cost $2$) and the third $v_2$-alive slot has only
$v_2$ alive (cost $3$), giving $0+2+2+3=7$; if $z_2\in S_4$ then $z_2$ contributes only
$v_4$ (cost $3$) and the remaining two $v_2$-alive slots cost $3$ each, giving
$3+2+3+3=11$. With $|S_4|=3$: if $z_4=z_2$ the three slots $\ne z_2$ have $v_2,v_4$ alive
(cost $2$ each) and $z_2$ is dead, giving $6$; if $z_4\ne z_2$ then $z_4$ has only $v_4$
alive (cost $3$), $z_2$ has only $v_4$ alive (cost $3$), and the two remaining slots have
$v_2,v_4$ (cost $2$ each), giving $3+3+2+2=10$. With $|S_4|=4$: $z_2$ has only $v_4$
(cost $3$) and the three others have $v_2,v_4$ (cost $2$), giving $3+2+2+2=9$. The minimum
is $6$, attained at $|S_4|=3$, $z_4=z_2$.

\smallskip
\noindent\emph{A3 (comp $3$ dead, conf full) and A4 (comp $3$ full, conf dead).}
In each case exactly two of the three components are nonzero constants, one of which is
alive on \emph{all four} slots. The slots in $S_4$ have two alive components (cost $2$) and
those outside $S_4$ have one alive component (cost $3$):
\[
  \text{cost}=2|S_4|+3(4-|S_4|)=12-|S_4|\ge 8 .
\]
The minima at $|S_4|=2,3,4$ are $10,9,8$ in each bucket.

\smallskip
\noindent\emph{A5 (comp $3$ full, conf co-point at $z_2$).}
Here $v_2=p(m+m(z_2))$ with $p\ne 0$, $v_3=c_3\ne 0$ constant, and $v_4$ as above.
With $|S_4|=2$: if $z_2\in S_4$ then $z_2$ has $v_3,v_4$ alive (cost $2$), the other slot of
$S_4$ has all three alive (cost $\ge 1$), and the two slots outside $S_4$ have $v_2,v_3$
alive (cost $2$ each), giving $\ge 2+1+2+2=7$; if $z_2\notin S_4$ then
$T=p^2(m+m(z_2))^2(c_3v_4)^{-1}$ has an injective numerator and a constant denominator on
$S_4$, so $T=1$ at most once, giving at most one cheap slot among the two three-alive slots
of $S_4$, hence $3+1+3+2=9$ (here $z_2$ is dead in $v_4$, contributing only $v_3$ at cost
$3$). With $|S_4|=3$: if $z_4=z_2$ the three slots $\ne z_2$ have all three alive
(cost $\ge 1$) and $z_2$ has only $v_3$ (cost $3$), giving $\ge 3+1+1+1=6$; if $z_4\ne z_2$
then $z_2$ has $v_3,v_4$ alive (cost $2$), $z_4$ has $v_2,v_3$ alive (cost $2$), and the two
remaining slots have all three alive (cost $\ge 1$ each), giving $\ge 2+2+1+1=6$.
With $|S_4|=4$: $v_4$ is two-valued and $(m+m(z_2))^2$ is injective, so $T=1$ at most once
per $v_4$-fibre pair, i.e.\ at most $2$ cheap slots total, giving $2+1+1+3=7$ (the dead-in-$v_4$
slot $z_2$ never occurs here since $|S_4|=4$; the bound is attained with one cheap slot per
pair and one three-alive slot at cost $3$). The minimum is $6$, attained at $|S_4|=3$.

\smallskip
\noindent\emph{A6 (comp $3$ full, conf full).}
Here $v_2,v_3$ are nonzero constants. With $|S_4|=2$: $v_4$ is constant on $S_4$, so the
cheapness equation $c_2^2=c_3v_4$ holds simultaneously on both slots of $S_4$ (all three
alive, cost $1$ each), while the two slots outside $S_4$ have $v_2,v_3$ alive (cost $2$),
giving $1+1+2+2=6$, which is tight. With $|S_4|=3$: $v_4$ is injective on $S_4$, so the
cheapness equation holds at most once, giving $\le 1$ cheap slot: $1+3+3+2=9$. With
$|S_4|=4$: $v_4$ is two-valued, so $\le 2$ cheap slots, giving $1+1+3+3=8$. The minimum is
$6$, attained at $|S_4|=2$.

\medskip
\noindent\textbf{Case $a\ne 0$ (comp $3$ a co-point at $z_3$, $v_3=a(m+m(z_3))$).}

\smallskip
\noindent\emph{B1 (conf dead).}
Here $v_2\equiv 0$ and $v_3$ is alive off $z_3$. If $b=1$ (so $|S_4|=3$): when $z_4=z_3$
the three slots $\ne z_3$ have $v_3,v_4$ alive (cost $2$) and $z_3$ is dead, giving $6$;
otherwise the cost is $3+3+2+2=10$. If $b\ne 1$ (double-paired): with $|S_4|=2$ and
$z_3\notin S_4$ the two $S_4$-slots have $v_3,v_4$ alive (cost $2$) and the third
$v_3$-alive slot has only $v_3$ (cost $3$), giving $0+2+2+3=7$ ($z_3\in S_4$ gives $11$ as
in A2); with $|S_4|=4$ the slot $z_3$ has only $v_4$ (cost $3$) and the three others have
$v_3,v_4$ (cost $2$), giving $3+2+2+2=9$. The minimum is $6$ (at $b=1$, $z_4=z_3$).

\smallskip
\noindent\emph{B2 (conf full).}
Here $v_2=c_2\ne 0$ constant and $v_3$ a co-point. If $b=1$: when $z_4=z_3$ the slot $z_3$
has only $v_2$ alive (cost $3$) and the three slots $\ne z_3$ have all three alive
(cost $\ge 1$), giving $\ge 3+1+1+1=6$; when $z_4\ne z_3$ the count is two slots with
$v_2,v_3$ or $v_2,v_4$ alive (cost $2$) and two all-three-alive slots (cost $\ge 1$),
giving $\ge 2+2+1+1=6$. If $b\ne 1$: with $|S_4|=2$ and $z_3\in S_4$, $z_3$ has $v_2,v_4$
alive (cost $2$), one all-alive slot (cost $\ge1$) and two $v_2,v_3$-alive slots (cost $2$),
giving $\ge 2+1+2+2=7$; with $z_3\notin S_4$, $T=c_2^2(v_3v_4)^{-1}$ has injective $v_3$ and
constant $v_4$ on $S_4$, so $\le 1$ cheap, giving $3+1+3+2=9$. With $|S_4|=4$: $v_3v_4$ is
injective on each $v_4$-fibre pair, so $\le 2$ cheap, giving $2+1+1+3=7$. The minimum
is $6$ (at $b=1$).

\smallskip
\noindent\emph{B3 (conf co-point at $z_2$, $b\ne 1$).}
Both $v_2$ and $v_3$ are co-points. If $z_2=z_3=:z$, the configuration and comp-$3$ values
are \emph{proportional} (both co-points of the line $\langle m\rangle$ at $z$), so
$T=(p^2/a)(m+m(z))\,v_4^{-1}$. With $|S_4|=2$ and $z\notin S_4$, $T$ has injective numerator
over a denominator constant on $S_4$, so $\le 1$ cheap, giving $0+1+3+2=6$ (tight); with
$z\in S_4$ the slot $z$ has only $v_4$ (cost $3$) and the survivors push the total to
$\ge 3+1+2+2=8$. With $|S_4|=4$, $\le 1$ cheap per $v_4$-fibre pair, giving $3+1+1+3=8$.
If $z_2\ne z_3$, then with $|S_4|=2$ each placement of $S_4$ relative to $\{z_2,z_3\}$ gives
$\ge 8$ (two $v_2,v_3$-alive slots at cost $4$ together with either two singly-alive slots
or a three-alive fibre join), and with $|S_4|=4$ the two slots off $\{z_2,z_3\}$ are alive in
all three components while $z_2,z_3$ each have two alive, giving $2+2+1+1=6$. The minimum
is $6$.

\smallskip
\noindent\emph{B4 (conf co-point, $b=1$).}
The deepest case. With $b=1$ the comp-$4$ direction is a bijection with dead slot $z_4$.
Suppose first $z_2=z_3=z_4=:z$. Then every slot $\ne z$ is alive in all three components,
and $z$ is free. Cheapness at a slot $s\ne z$ reads
\[
  p^2\bigl(m(s)+m(z)\bigr)^2=a\bigl(m(s)+m(z)\bigr)\,v_4(s),
  \qquad\text{i.e.}\qquad
  \bigl(m(s)+m(z)\bigr)\bigl(k_4(s)+k_4(z)\bigr)^{-1}=a\,p^{-2},
\]
a level condition on the chord slope of the quadruple $(m,k_4)=(m,\,m+\omega\theta)$.
By the chord-slope lemma (Lemma~\ref{lem:chord}) the slope is constant on a triple of slots
iff the corresponding points of $AG(2,\F4)$ are collinear; and by the hyperbolic-quadruple
lemma (Lemma~\ref{lem:hyperbolic}) the quadruple $\{(m(s),k_4(s))\}=H_{\omega^2}$ has no
three of its four points collinear. Hence the level condition is met on at most one slot:
$\le 1$ cheap, giving $0+1+3+3=7$. The remaining alignments of $\{z_2,z_3,z_4\}$ all give
$\ge 7$:
\begin{itemize}[leftmargin=2em]
  \item exactly two of $z_2,z_3,z_4$ coincide: the doubled slot is at most singly alive
        (cost $\ge 3$ if its $v_0$-free cost is $3$ -- e.g.\ $z_2=z_3\ne z_4$ leaves the
        shared slot with only $v_4$ at cost $3$ -- otherwise the third, dead slot costs $2$),
        and at most two three-alive slots remain, forcing $\ge 7$ in each of the three
        patterns;
  \item all of $z_2,z_3,z_4$ distinct: three slots carry two alive components (cost $2$
        each) and one slot is three-alive (cost $\ge 1$), giving $2+2+2+1=7$.
\end{itemize}
The minimum is $7$ across this bucket.

\medskip
Every bucket has minimum $\ge 6$, so $S(a,b)\ge 6$ for all $16$ pairs $(a,b)$, and in fact
$S(a,b)=6$ for every $(a,b)$ since each row above attains a configuration of cost $6$ for
the relevant parameters. By Lemma~\ref{lem:stdform} the four weight-$24$ block tables
inherit the value $6$ everywhere. Since the per-block linked minimum at a spine cell is the
sum of the two block costs and the $v_0$-free cost lower-bounds each, every weight-$24$
spine cell has linked value $\ge 6+6=12$, which is the assertion (M-im) on the six
weight-$24$ Smith classes.
\end{proof}

The bucket walk is summarized in Table~\ref{tab:wt24buckets}; each entry is the minimum of
the displayed slot-cost computation over the indicated fibre type, and the reader may
re-derive any single entry directly from the slot-cost rule.

\begin{table}[h]
\centering
\begin{tabular}{@{}lc@{\hspace{3em}}lc@{}}
\toprule
bucket & min & bucket & min \\
\midrule
A1\ \ $|S_4|=2/3/4$ & $6/9/12$ & B1\ \ $b=1$ & $6$ \\
A2\ \ $|S_4|=2/3/4$ & $7/6/9$  & B1\ \ $b\ne1,\ |S_4|=2/4$ & $7/9$ \\
A3\ \ $|S_4|=2/3/4$ & $10/9/8$ & B2\ \ $b=1$ & $6$ \\
A4\ \ $|S_4|=2/3/4$ & $10/9/8$ & B2\ \ $b\ne1,\ |S_4|=2/4$ & $7/7$ \\
A5\ \ $|S_4|=2/3/4$ & $7/6/7$  & B3\ \ $z_2=z_3,\ |S_4|=2/4$ & $6/8$ \\
A6\ \ $|S_4|=2/3/4$ & $6/9/8$  & B3\ \ $z_2\ne z_3,\ |S_4|=2/4$ & $8/6$ \\
        &          & B4\ \ $z_2{=}z_3{=}z_4 / z_2{=}z_3 / z_2{=}z_4 / z_3{=}z_4 / \text{distinct}$ & $7/7/9/9/7$ \\
\bottomrule
\end{tabular}
\caption{Bucket minima for the standard-form problem $S(a,b)$ (the $v_0$-free block costs).
Every entry is $\ge 6$, and the row minimum is $6$ for each $(a,b)$; hence $S(a,b)=6$ for
all $16$ pairs.}
\label{tab:wt24buckets}
\end{table}
\section{The weight-16 and weight-18 orbits: locus rules and the per-cell floor}
\label{sec:wt1618}

On the three light orbits --- which we label \emph{wt-16}, \emph{wt-18a}, and
\emph{wt-18b} after the cell value of their generic spine cell --- the unlinked block
floors are only $3$ to $5$. It is the linkage carried by the shared slot data $(V_0,\gamma)$
that lifts the per-cell floor to $10$. Working in the slot frame, we produce, per orbit and
per spine cell, enough of the function
\[
  (V_0,\gamma)\ \longmapsto\ (\min_L,\ \min_R)
\]
to certify the two requirements isolated by the achiever-structure lemma
(Lemma~\ref{lem:achiever}): the cost-$8$ kill ($\min_L+\min_R\ge 10$ for every
$(V_0,\gamma)$) and the explicit loci where $\min_L+\min_R=10$, together with their
argmins.

Throughout, slots are the four elements of $\Z_2^2=\{e,x,y,xy\}$; a component datum is a
function from slots to $\F4$, written as a $4$-tuple over $(e,x,y,xy)$. We carry the slot
labelings $m=(\omega^2,\omega,1,0)$ and $m'=(\omega^2,1,\omega,0)=m\circ(x\leftrightarrow y)$,
the constants $\theta=(1,0,1,0)$ and $\tilde\theta=(1,1,0,0)$, the confining line of a block,
the comp-$3$ direction $k_3=\kappa_{3}+a_3 m$ (resp.\ $m'$ on $R$), and the comp-$4$
direction with its tied constant $d_4=\omega\gamma+e$ (the component-$4$ tie of the slot
frame). The
``alive'' count of a slot is the number of nonzero values among its three constrained
components; the per-slot cost rule and the slot-parity identity $\text{cost}\equiv|V_0|
\pmod 2$ are those of the slot-cost lemma (Lemma~\ref{lem:slotcost}), with the cheapness
test $T=1$, $T_L=v_2^2(v_3v_4)^{-1}$ on $L$ and $T_R=v_4(v_1v_3)^{-1}$ on $R$.

\subsection{The locus rules}\label{ssec:locus-rules}

We record five rules that turn a per-block slot-cost target into a finite list of admissible
configurations. Each is a direct consequence of the slot-cost lemma
(Lemma~\ref{lem:slotcost}), the pair-ratio lemma for affine pencils (Lemma~\ref{lem:pencil}),
and the chord-slope/hyperbolic facts (Lemma~\ref{lem:chord}, Lemma~\ref{lem:hyperbolic}).

\begin{description}
\item[(R1) Zero slot.] A slot costs $0$ if and only if $v_0=0$ there and all three
  constrained components vanish at that slot.

\item[(R2) Dead-pair rigidity.] If two slots $\{s,s'\}$ both cost $0$, then the confining
  component is identically $0$ (a nonzero point of the confining line vanishes at most at one
  slot), and $\{s,s'\}$ lies inside a single fibre of \emph{both} the comp-$3$ and the
  comp-$4$ direction. By the fibre tables of the pair-ratio lemma (Lemma~\ref{lem:pencil}),
  this pins the spine and the surviving constants --- hence $\gamma$ --- up to the listed
  fibre choices. Three zero slots in addition force $v_3\equiv 0$ on them, i.e.\ a fibre of
  size $\ge 3$ in the comp-$3$ direction; such a fibre exists only where comp $3$ is
  confining-parallel and the spine annihilates it (the wt-18a row $a_3=1$), with a $3{+}1$
  fibre of $k_4$ carrying the remaining cost.

\item[(R3) $\delta$-slot.] A slot with $v_0=1$ costs $1$ if and only if all three components
  are alive there and $T=1$. The confining scale $p$ enters $T$ quadratically on $L$ and
  linearly on $R$; matching two $\delta$-slots (or one $\delta$-slot and one cheap
  three-alive slot) against a single $p$ is one consistency equation in $\F4$, solved or
  refuted by inspection.

\item[(R4) Cost-2 slot.] A slot with $v_0=0$ costs $2$ if and only if exactly two components
  are alive there, or all three are alive with $T=1$.

\item[(R5) Shapes by parity.] By slot parity the admissible per-block slot-cost multisets at
  total cost $3,4,5$ are exactly:
  \[
  \begin{aligned}
  \text{cost }3:\ & \{3,0,0,0\},\ \{1,2,0,0\}\ (|V_0|{=}1),\ \{1,1,1,0\}\ (|V_0|{=}3);\\
  \text{cost }4:\ & \{4,0,0,0\},\ \{2,2,0,0\}\ (|V_0|{=}0),\ \{1,1,2,0\},\ \{1,3,0,0\}\ (|V_0|{=}2),\ \{1,1,1,1\}\ (|V_0|{=}4);\\
  \text{cost }5:\ & \{1,2,2,0\},\ \{1,4,0,0\},\ \{3,2,0,0\}\ (|V_0|{=}1),\ \{1,1,3,0\},\ \{1,1,1,2\}\ (|V_0|{=}3).
  \end{aligned}
  \]
\end{description}

Each locus-table row below is obtained by the same finite procedure: choose an admissible
shape by (R5), apply (R1)--(R4) to pin the configuration data, and test the resulting
consistency equations in $\F4$. We exhibit one cell per orbit in full; the remaining cells
are settled by repeating this procedure, and the complete loci are collected in the tables of
\S\ref{ssec:wt1618-tables}, where every entry is a direct hand computation.

\subsection{The weight-16 orbit}\label{ssec:wt16}

On the $L$ block the confining line is $\langle m\rangle$, and
\[
  v_3=a_3 m+\omega\theta+c_3,\qquad v_4=a_4 m+\omega\theta+c_4,\qquad d_{4L}=\omega\gamma+\omega .
\]
On the $R$ block, after applying Frobenius to the comp-$4$ values (a cost-preserving move),
the confining line is $\langle m'\rangle$ with
\[
  v_3=a_3 m'+\theta+c_3,\qquad \tilde v_4=a_4^2 m'+\theta+c_4 .
\]
The comp-$3$ and comp-$4$ directions are double-paired for a spine value in
$\{0,\omega,\omega^2\}$ and bijective at $1$. On $L$ the pairings are $\{e,y\mid x,xy\}$ at
$0$, $\{e,x\mid y,xy\}$ at $\omega$, and $\{e,xy\mid x,y\}$ at $\omega^2$; on $R$ (whose
offset is $\theta$ against $m'$) the pairings at $\omega$ and $\omega^2$ are exchanged. The
sixteen spine cells are indexed by $(a_3,a_4)\in\F4^2$.

\paragraph{Worked cell $(\omega^2,\omega^2)$, the $(4,6)$ locus.}
Consider the cost-$4$ $L$-shapes. The shape $\{2,2,0,0\}$ requires a dead pair; by (R2) the
confining component is $\equiv 0$ and the dead pair lies in a common fibre of
$k_3=\omega^2 m+\omega\theta$ and $k_4=\omega^2 m+\omega\theta$. These directions are equal,
with pairing $\{e,xy\mid x,y\}$ (pair-ratio $\omega^2$ at spine value $\omega^2$), and both
fibres are available:
\begin{itemize}
\item dead pair $\{e,xy\}$ pins $c_3=c_4=k(e)$ and $d_{4L}=k_4(e)=\omega^2 m(e)+\omega\theta(e)
  =\omega+\omega=0$, so $\omega\gamma+\omega=0$ gives $\gamma=1$;
\item dead pair $\{x,y\}$ pins $d_{4L}=k_4(x)=1$, hence $\gamma=\omega$.
\end{itemize}
In each case the two live slots carry $v_3=v_4=$ the pair-gap value, with two components
alive each, contributing cost $2+2$ at $V_0=0000$. Thus $\min_L=4$ at
$(V_0,\gamma)\in\{(0000,1),(0000,\omega)\}$, with a unique argmin. The all-$\delta$ shape
$\{1,1,1,1\}$ ($|V_0|=4$) also solves here: since $k_3=k_4$, the $\delta$-condition $T\equiv1$
reads $p^2(m\text{-line})^2=v_3 v_4=v_3^2$, one equation per slot, and these are consistent
across all four slots exactly at $\gamma\in\{0,\omega^2\}$ (with $V_0=1111$). The $|V_0|=2$
shapes and $\{4,0,0,0\}$ fail: with $k_3=k_4$ the equations of (R3) are overdetermined. The
$R$ block at each of these four $(V_0,\gamma)$ has $\min_R=6$. Hence this cell contributes
$(4,6)$ achievers at the four loci $1111@0$, $0000@1$, $0000@\omega$, $1111@\omega^2$, each
with $|\mathrm{Argmin}_R|=2$, matching the table.

\paragraph{Worked cell $(\omega,1)$, a $(5,5)$ locus.}
Take $V_0=\delta_{xy}$ and $\gamma=0$, so $d_{4L}=\omega$. Consider the shape
$\{1,2,2,0\}$. The zero slot $s_0$ must satisfy $k_4(s_0)=d_{4L}$; with
$k_4=m+\omega\theta=(1,\omega,\omega^2,0)$ (bijective, $a_4=1$) this forces $s_0=x$. By (R1)
the confining component is the co-point $p\,(m+m(x))$. The comp-$3$ direction
$k_3=\omega m+\omega\theta=(\omega^2,\omega^2,0,0)$ is double-paired $\{e,x\mid y,xy\}$, so
$v_3$ vanishes on the whole pair $\{e,x\}$; consequently slot $e$ has exactly the confining
and the comp-$4$ components alive --- cost $2$, as required. Slots $y$ and $xy$ have all three
components alive; the $\delta$-slot is $xy$ (with $v_0=1$), and slot $y$ must satisfy $T=1$ by
(R4). Each yields one equation for $p^2$:
\[
  \text{at }xy:\quad p^2\cdot\omega^2=v_3(xy)\,v_4(xy)=\omega^2\cdot\omega,\quad\text{so }p^2=\omega;
  \qquad
  \text{at }y:\quad p^2\cdot(\omega^2)^2=\omega^2\cdot 1,\quad\text{so }p^2=\omega.
\]
The two are consistent, giving $p=\omega^2$. The configuration exists and is unique: a
$(5,5)$ locus with singleton argmins. The other three loci $(V_0,\gamma)=(\delta_s,\cdot)$ are
its images under the slot translations: the stabilizer of the orbit has order $4$ and acts by
the three nonzero slot translations, which permute the $\delta$-positions and shift $\gamma$
accordingly.

\paragraph{The remaining cells.}
Each of the remaining cells is settled by the same finite procedure; the complete loci are
collected in Table~\ref{tab:wt16}. We record the qualitative outcome that the cost-$8$ kill
depends on.
\begin{itemize}
\item \emph{Floor-$10$ cells} $(\omega,1)$, $(\omega,\omega^2)$, $(\omega^2,1)$,
  $(\omega^2,\omega^2)$: each contributes four $(5,5)$ loci with singleton argmins and four
  $(6,4)$- or $(4,6)$-loci with two argmin pairs each, for $4+8=12$ achievers per cell and
  $48$ in total.
\item \emph{All other cells} realize the $(4,4)$ kill, i.e.\ the disjointness of the $L4$ and
  $R4$ loci read off Table~\ref{tab:wt16}: at every such cell one side is empty, or the
  $V_0$-sets are disjoint (e.g.\ at $(\omega,1)$, $V_0(L4)\subseteq\{0011,1100\}$ versus
  $V_0(R4)\subseteq\{1001,0110\}$), or the $V_0$'s coincide and the $\gamma$-sets are
  complementary (e.g.\ at $(0,0)$, $L4$ at $(0000,\{0,1\})\cup(1111,\{\omega,\omega^2\})$ and
  $R4$ at $(1111,\{0,1\})\cup(0000,\{\omega,\omega^2\})$).
\item \emph{$L3=R3=\varnothing$ at every cell}, fully derived as follows. The $|V_0|=1$ shapes
  $\{1,2,0,0\}$ and $\{3,0,0,0\}$ both need $\ge 2$ zero slots, killing the confining
  component by (R2): the former then has no alive confining component at its $\delta$-slot,
  and the latter needs $v_3\equiv 0$ on three slots, impossible since $\kappa_{3L}=\omega\theta$
  is non-constant and $k_3$ has no fibre of size $\ge 3$. The $|V_0|=3$ shape $\{1,1,1,0\}$
  (one zero slot $s_0$, three $\delta$-slots) dies in three sub-cases: at $a_3\ne 1$ the
  comp-$3$ direction is double-paired, so $v_3(s_0)=0$ kills $v_3$ on $s_0$'s partner slot,
  contradicting that slot being a $\delta$-slot; at $(1,a_4)$ with $a_4\ne 1$ the same
  argument applies to comp $4$; and at $(1,1)$, where $k_3=k_4=m+\omega\theta$, the triple-$\delta$
  condition says $(m+m(s_0))\cdot(k_3+k_3(s_0))^{-1}$ is constant on the three slots
  $\ne s_0$, i.e.\ three points of the hyperbolic quadruple $(m,m+\omega\theta)=H_{\omega^2}$
  are collinear, contradicting Lemma~\ref{lem:hyperbolic}.
\end{itemize}
With slot parity, every $(V_0,\gamma)$ satisfies $\min_L+\min_R\ge 10$, so $m(\text{cell})\ge
10$ for all $16$ cells, with equality exactly at the four floor-$10$ cells.

\subsection{The weight-18a orbit}\label{ssec:wt18a}

On $L$, $v_3=(1+a_3)m+c_3$ --- here $\kappa_{3L}=m$ \emph{is} the labeling, so comp $3$ is
confining-parallel, dead or full at $a_3=1$ and a co-point otherwise --- while $v_4$ has
direction $k_4=(0,\omega,\omega^2,0)+a_4 m$ with fibre types $2{+}1{+}1$, $2{+}1{+}1$,
$3{+}1$, $2{+}1{+}1$ over $a_4=0,1,\omega,\omega^2$, and $d_{4L}=\omega\gamma+\omega^2$. On
$R$, $v_3=(1+a_3)m'$ and $\tilde v_4$ has direction $\omega(1,1,1,0)+a_4^2 m'$ (type $3{+}1$
at $a_4=0$, else $2{+}1{+}1$). The translation stabilizer of the orbit has order $3$ and acts
on the spine by affine maps of the form $a_3\mapsto\omega(1+a_3)+1$ that fix $a_4$: the cells
$(a_3,a_4)$ with $1+a_3\ne 0$ form orbits of three within each $a_4$-column, and the row
$a_3=1$ is fixed. There are therefore six cell classes: one rep $(\tilde a_3\ne 0,a_4)$ for
each $a_4\in\F4$, together with $(1,0)$ and $(1,\omega)$. (The cells $(1,1)$ and
$(1,\omega^2)$ have cell value $12$ with empty low loci.)

\paragraph{Worked class rep $(0,1)$, the $|V_0|=3$ loci.}
The shape $\{1,1,1,2\}$ has its (R3) $\delta$-equations solvable at exactly
$(V_0,\gamma)\in\{(1011,1),(1101,\omega^2)\}$ on the $(5,5)$ side, together with the $(6,4)$
locus $(1100,\omega^2)$. The class transports to $(\omega,1)$ and $(\omega^2,1)$ verbatim, as
the stabilizer fixes $\gamma$ and $a_4$ and permutes nothing else in the table. The isolated
row $a_3=1$ kills comp $3$ entirely (dead or full): at $(1,0)$ the surviving loci are $(6,4)$
at $(0000,0)$ and $(7,3)$ at $(0001,0)$; at $(1,\omega)$ one finds the mirror $(3,7)/(4,6)$.
The complete count over the fourteen floor-$10$ cells is $48$ achievers: $12$ per $a_4\in
\{1,\omega^2\}$ column-orbit, $6$ per $a_4\in\{0,\omega\}$ column-orbit, and $6+6$ at the two
fixed cells.

\paragraph{The cost-8 kill.}
All $L4$, $R4$, $L3$, $R3$ loci are singletons (Table~\ref{tab:wt18a}) and pairwise disjoint
in $(V_0,\gamma)$ at every cell. The only $L3/R3$ sites occur at $(1,\omega)$ and $(1,0)$,
where the partner block's minimum at that $(V_0,\gamma)$ is $7$.

\subsection{The weight-18b orbit}\label{ssec:wt18b}

On $L$, $\kappa_{3L}=(1,\omega,\omega^2,0)$ is a bijective offset, so
$k_3=\kappa_{3L}+a_3 m$ has fibre types $1{+}1{+}1{+}1$, $2{+}2$, $2{+}2$, $2{+}2$ over
$a_3=0,1,\omega,\omega^2$; $\kappa_{4L}=(\omega,\omega,1,0)$ gives $k_4$ fibre types
$2{+}1{+}1$, $3{+}1$, $2{+}1{+}1$, $2{+}1{+}1$; and $d_{4L}=\omega\gamma+1$. On $R$,
$\kappa_{3R}=\omega^2 m$ (a fixed non-affine offset in the $m'$-frame) and
$\tilde\kappa_{4R}=(\omega^2,\omega,\omega^2,0)$. The translation stabilizer is trivial, but
the orbit carries a swap-type symmetry: composing translation with the block swap fixes the
orbit, exchanges the two blocks, and pairs the cells. Each locus table at $(a_3,a_4)$ thus
mirrors a partner cell's with $L$ and $R$ exchanged --- visible in Table~\ref{tab:wt18b},
e.g.\ $(0,1)$ carries $(3,7)+(4,6)$ while $(0,\omega^2)$ carries $(7,3)+(6,4)$ with the same
$V_0$'s. There are $12$ floor-$10$ cells and $22$ achievers; the two $L3/R3$ sites
($(0,1)$ and $(0,\omega^2)$, both at $V_0=0111$, the unique $|V_0|=3$ cost-$3$ shape
$\{1,1,1,0\}$ the frame admits) have partner minima $7$.

\subsection{The locus tables and the C-table floor}\label{ssec:wt1618-tables}

In the tables below $V_0$ is written as a bit-string over the slots $(e,x,y,xy)$ and
$\gamma\in\F4$; a locus is recorded as $V_0@\gamma$. Each row is the outcome of the finite
procedure of \S\ref{ssec:locus-rules}, and is small enough to be checked by hand.

\begin{table}[h]
\centering
\caption{Weight-16 floor-$10$ cells and their achiever loci.}
\label{tab:wt16}
\begin{tabular}{lll}
\toprule
cell & $(5,5)$ loci (argmins $1\times1$) & $(4,6)/(6,4)$ loci (argmins $1\times2$ / $2\times1$)\\
\midrule
$(\omega,1)$        & $0001@0,\ 0100@1,\ 0010@\omega,\ 1000@\omega^2$ & $(6,4)$: $1001@0,\ 0110@1,\ 0110@\omega,\ 1001@\omega^2$\\
$(\omega,\omega^2)$ & $0111@0,\ 1110@0,\ 1011@\omega^2,\ 1101@\omega^2$ & $(6,4)$: $1111@0,\ 0000@1,\ 0000@\omega,\ 1111@\omega^2$\\
$(\omega^2,1)$      & $1000@0,\ 0010@1,\ 0100@\omega,\ 0001@\omega^2$ & $(4,6)$: $1001@0,\ 0110@1,\ 0110@\omega,\ 1001@\omega^2$\\
$(\omega^2,\omega^2)$ & $1011@0,\ 1101@0,\ 0111@\omega^2,\ 1110@\omega^2$ & $(4,6)$: $1111@0,\ 0000@1,\ 0000@\omega,\ 1111@\omega^2$\\
\bottomrule
\end{tabular}
\end{table}

\noindent $12$ achievers per cell, $48$ in total. At every other cell $L4\cap R4=\varnothing$
and $L3=R3=\varnothing$, so those cells have value $\ge 12$.

\begin{table}[h]
\centering
\caption{Weight-18a cell classes (under the order-$3$ stabilizer; rows $\tilde a_3=1+a_3\ne0$
transport, the row $a_3=1$ is fixed) and their loci.}
\label{tab:wt18a}
\begin{tabular}{ll}
\toprule
class & loci\\
\midrule
$(\tilde a_3\ne 0,\,0)$        & $(4,6)$: $0110@\omega$\\
$(\tilde a_3\ne 0,\,1)$        & $(5,5)$: $1011@1,\ 1101@\omega^2$;\quad $(6,4)$: $1100@\omega^2$\\
$(\tilde a_3\ne 0,\,\omega)$   & $(6,4)$: $0110@\omega^2$\\
$(\tilde a_3\ne 0,\,\omega^2)$ & $(4,6)$: $1100@\omega$;\quad $(5,5)$: $1011@0,\ 1101@\omega$\\
$(1,\,0)$                      & $(6,4)$: $0000@0$;\quad $(7,3)$: $0001@0$\\
$(1,\,\omega)$                 & $(3,7)$: $0001@1$;\quad $(4,6)$: $0000@1$\\
\bottomrule
\end{tabular}
\end{table}

\noindent Achiever counts: $2$ per $(\tilde a_3\ne 0,0)$- and $(\tilde a_3\ne 0,\omega)$-cell,
$4$ per $(\tilde a_3\ne 0,1)$- and $(\tilde a_3\ne 0,\omega^2)$-cell, and $6$ at each fixed
cell, giving $3\cdot(2+4+2+4)+6+6=48$. The cells $(1,1)$ and $(1,\omega^2)$ have value $12$
with empty loci.

\begin{table}[h]
\centering
\caption{Weight-18b loci, paired by the swap symmetry (each left row mirrors its right
partner with the $(c_L,c_R)$ split exchanged).}
\label{tab:wt18b}
\begin{tabular}{llcll}
\toprule
cell & loci & & swap-partner & loci\\
\midrule
$(0,1)$        & $(3,7)$: $0111@0$;\ $(4,6)$: $0101@0$        & & $(0,\omega^2)$ & $(6,4)$: $0101@\omega$;\ $(7,3)$: $0111@\omega$\\
$(1,\omega)$   & $(4,6)$: $0000@1$                            & & $(1,0)$        & $(6,4)$: $0000@\omega^2$\\
$(1,\omega^2)$ & $(4,6)$: $1010@1$;\ $(5,5)$: $0010@1$        & & $(1,1)$        & $(5,5)$: $0010@\omega^2$;\ $(6,4)$: $1010@\omega^2$\\
$(\omega,0)$   & $(4,6)$: $1001@0$;\ $(5,5)$: $1110@1$;\ $(6,4)$: $1100@1$ & & $(\omega^2,\omega)$ & $(4,6)$: $1100@\omega^2$;\ $(5,5)$: $1110@\omega^2$;\ $(6,4)$: $1001@\omega$\\
$(\omega,1)$   & $(5,5)$: $1011@\omega$                       & & $(\omega^2,\omega^2)$ & $(5,5)$: $1011@0$\\
$(\omega^2,0)$ & $(4,6)$: $0000@\omega$;\ $(5,5)$: $0001@\omega$ & & $(\omega,\omega)$ & $(5,5)$: $0001@0$;\ $(6,4)$: $0000@0$\\
\bottomrule
\end{tabular}
\end{table}

\noindent Achiever count $2+2+1+1+2+2+3+3+1+1+2+2=22$.

\begin{proposition}[C-table floor]\label{prop:ctablefloor}
Every Smith-coset element has weight $\ge 12$ on the wt-24 orbits and $\ge 10$ on the
wt-16, wt-18a, and wt-18b orbits. On each of the latter, for every $(V_0,\gamma)$ the per-block
minima satisfy $\min_L+\min_R\ge 10$, with equality exactly on the loci listed in
Tables~\ref{tab:wt16}--\ref{tab:wt18b}. The cost-$10$ configurations are therefore exactly the
$48+48+22=118$ achievers obtained from those loci.
\end{proposition}

\begin{proof}
On the wt-24 orbits the bound is $S(a,b)\ge 6$ for every block at every cell
(Proposition~\ref{prop:wt24}); since the confining-free cost lower-bounds every fixed-$V_0$
cost, each wt-24 spine cell has linked value $\ge 6+6=12$.

For each light orbit the work above produces, cell by cell, the two facts the
achiever-structure lemma (Lemma~\ref{lem:achiever}) requires. First, $\min_L+\min_R\ge 10$ for
every $(V_0,\gamma)$: by slot parity (Lemma~\ref{lem:slotcost}) only the splits $(3,3)$, $(4,4)$,
$(3,5)$ and $(5,3)$ can violate this, and these are excluded cell by cell --- the $(4,4)$ kill by
the disjointness of the $L4$ and $R4$ loci, and the $(3,3)/(3,5)/(5,3)$ kills because every
cost-$3$ locus (an $L3$ or $R3$ site) pairs with a partner-block minimum $\ge 7$ at the same $(V_0,\gamma)$
(\S\ref{ssec:wt16}--\ref{ssec:wt18b}). Second, the loci where $\min_L+\min_R=10$ are exactly
those tabulated, with the stated argmin multiplicities. By Lemma~\ref{lem:achiever} the
cost-$10$ configurations are precisely
\[
  \bigcup_{(V_0,\gamma):\,\min_L+\min_R=10}\ \mathrm{Argmin}_L(V_0,\gamma)\times\mathrm{Argmin}_R(V_0,\gamma),
\]
which the tables enumerate as $48$ (wt-16) $+48$ (wt-18a) $+22$ (wt-18b) $=118$ achievers.
\end{proof}

\section{The $\rho$-link kills: no weight-$10$ element}\label{sec:rho-kills}

By the achiever-structure lemma (Lemma~\ref{lem:achiever}), a weight-$10$ Smith-coset element
must realize one of the $118$ achievers of Proposition~\ref{prop:ctablefloor}, with its free
sides in the per-slot argmin sets and \emph{both} $\rho$-links satisfied:
\[
  V_{1R}=\rho_1\cdot V_{1L},\qquad V_{2L}=\rho_2\cdot V_{2R}
\]
(Lemma~\ref{lem:offvanish-statement}). For an achiever the values $V_{1R}$ and $V_{2L}$ are already part
of the configuration; the links then demand $V_{1L}\in\rho_1^{-1}(V_{1R})$, a coset of
$\ker\rho_1=\F4\hA_1\oplus\F4\,XY$ (sixteen elements, nonempty since $V_{1R}\in\img\rho_1$),
meeting the product $\prod_s\mathrm{Argmin}_1(s)\subseteq\F4$ of the four per-slot comp-$1$
argmin sets, and the mirror statement for $V_{2R}$. The per-slot argmin sets are read off the
value table: at a slot of cost $c$, $\mathrm{Argmin}_1(s)=\{v_1: \text{the layer with values }
(v_0;v_1,v_2,v_3,v_4)\text{ has weight }c\}$ --- almost always a singleton, since a
$\delta$-point pins all five values and a weight-$2$ slot pins them up to its single dead
component.

\begin{proposition}[the $\rho$-link kill]\label{prop:kill}
Each of the $118$ weight-$10$ achievers of Proposition~\ref{prop:ctablefloor} violates at
least one $\rho$-link. Precisely, $116$ achievers violate both links and the remaining $2$
(the wt-18b $(5,5)$ achievers at cells $(\omega,\omega)$ and $(\omega^2,0)$) violate exactly
one. Consequently no Smith-coset element has weight $10$.
\end{proposition}

\begin{proof}
The check is, per achiever, a single convolution in $\F4[\Z_2^2]$ followed by a slot-by-slot
comparison. We carry out the worked head of the wt-16 $(5,5)$ family in full and then describe
the family table; each remaining check is one such evaluation.

\paragraph{Worked kill: the wt-16 $(5,5)$ head.}
At the cell $(\omega,1)$, with $V_0=\delta_{xy}=0001$ and $\gamma=0$ (the cell of the worked
$(5,5)$ locus of \S\ref{ssec:wt16}), the achiever's $L$ data is
\[
  V_{2L}=(\omega^2,0,\omega,1),\quad V_{3L}=(0,0,\omega^2,\omega^2),\quad V_{4L}=(\omega^2,0,1,\omega),
\]
with slot costs $[2,0,2,1]$. The per-slot comp-$1$ argmin sets are the singletons
$(\{1\},\{0\},\{0\},\{\omega^2\})$: at the zero slot the empty layer forces $v_1=0$; at the
$\delta$-slot the $\delta$-point pins $v_1$; and at the two weight-$2$ slots the one dead
component of the pair is not comp $1$. Hence link $1$ holds if and only if the single
convolution $\rho_1\cdot(1,0,0,\omega^2)$ equals $V_{1R}=(\omega^2,\omega,0,1)$. Computing in
the $X,Y$ basis, $(1,0,0,\omega^2)=\omega\cdot 1+\omega^2 X+\omega^2 Y+\omega^2 XY$, and
\[
  (X+\omega Y+\omega^2 XY)\cdot(\omega+\omega^2 X+\omega^2 Y+\omega^2 XY)=\omega X+\omega^2 Y+\omega^2 XY,
\]
whose slot values are $(\omega,1,0,\omega^2)\ne(\omega^2,\omega,0,1)=V_{1R}$: link $1$ fails.
The mirror computation against $\rho_2$ kills link $2$ as well, so this achiever violates both
links.

\paragraph{The remaining achievers.}
The same one-convolution check settles every achiever: for each, one forms $\rho_i$ times the
unique product of the comp-$i$ argmin singletons and compares the four slot values against the
configuration's $V_{iR}$ (resp.\ $V_{iL}$). Equivalently, the product of the four argmin sets
and the sixteen-element link coset are disjoint, which is a single $\F4$ evaluation per slot.
Transporting the worked head by the orbit symmetries of \S\ref{sec:wt1618} sweeps each family.
The outcome is recorded in Table~\ref{tab:kills}: $116$ achievers fail both links, and the two
exceptions noted in the statement fail exactly one. In every case at least one link fails, so
none of the $118$ achievers is link-compatible.
\end{proof}

\begin{table}[h]
\centering
\caption{The $118$ weight-$10$ achievers and their $\rho$-link kills.}
\label{tab:kills}
\begin{tabular}{lll}
\toprule
family & achievers & links failed\\
\midrule
wt-16   & $48$ & both\\
wt-18a  & $48$ & both\\
wt-18b (non-exceptional) & $20$ & both\\
wt-18b $(5,5)$ at $(\omega,\omega),\,(\omega^2,0)$ & $2$  & exactly one\\
\midrule
total & $118$ & all fail $\ge 1$\\
\bottomrule
\end{tabular}
\end{table}

\begin{proposition}\label{thm:Mim}
Every base $1$-cycle lying in a nonzero Smith class $\img\cmap$ has weight $\ge 12$.
\end{proposition}

\begin{proof}
By the standard-form walk (Proposition~\ref{prop:wt24}) no Smith-coset element has cost
$\le 9$, the wt-24 orbits being floored at $12$ and the wt-16/18 orbits at $10$
(Proposition~\ref{prop:ctablefloor}). By the $\rho$-link kill (Proposition~\ref{prop:kill}) no
link-compatible cost-$10$ configuration exists. Finally, every Smith-coset element has even
weight (the evenness of the slot frame). Combining: a Smith-coset element of weight $<12$ would
have weight $10$ and would realize one of the $118$ achievers with both links satisfied,
contradicting Proposition~\ref{prop:kill}; hence every such element has weight $\ge 12$.

This bounds the five orbit representatives. Transport across the Smith classes (the slot-frame
symmetry that carries one nonzero class to another preserving weight) extends the bound from
the representatives to all $63$ nonzero classes of $\img\cmap$. Therefore every base $1$-cycle
in a nonzero Smith class has weight $\ge 12$.
\end{proof}
\section{Theorem D: the distance}\label{sec:thmD}

We now assemble the pieces into the distance computation. All weights are taken on
the $Z$ side; the $X$ side is identical by the Duality Lemma (Lemma~\ref{lem:duality}).

\begin{thmD}
$d(\mathrm{gross})=12$.
\end{thmD}

\begin{proof}
\emph{Lower bound.} Let $v=(v_L,v_R)$ be a nontrivial $Z$-logical of the gross code,
i.e.\ a cover cycle whose class in $H_1(\mathrm{gross})$ is nonzero. Write
$p(v)=v_0+v_1$ for its sheet-sum projection to the base, a chain map with
$\wt v\ge\wt{p(v)}$. We split on the class of $p(v)$ in $H_1(\mathrm{base})$.

If $[p(v)]=0$, then $v$ lies in the kernel of the induced projection
$p_\ast\colon H_1(\mathrm{gross})\to H_1(\mathrm{base})$, so its class belongs to the
dangerous sector. Theorem~C applies directly and gives $\wt v\ge 12$.

If $[p(v)]\ne 0$, then by the homotopy theorem~(R) we have
$[p(v)]\in\img\cmap$; combined with $[p(v)]\ne 0$ this places $[p(v)]$ in
$\img\cmap\smallsetminus 0$. Thus $p(v)$ is a base $1$-cycle in a nonzero Smith class,
and Proposition~\ref{thm:Mim} gives $\wt{p(v)}\ge 12$. Hence
$\wt v\ge\wt{p(v)}\ge 12$.

In either case $\wt v\ge 12$, so $d_Z(\mathrm{gross})\ge 12$. By the Duality Lemma
(Lemma~\ref{lem:duality}), $d_X(\mathrm{gross})=d_Z(\mathrm{gross})$, and therefore
$d(\mathrm{gross})\ge 12$.

\emph{Upper bound and tightness.} Let
$z^\ast=1+y+y^2+y^5+x^3+x^3y^4\in\Ann(A)$ be the weight-$6$ annihilator element of
Corollary~A$'$ (\ref{cor:Aprime}), and let $u^\ast=(z^\ast,0)$ be the associated
weight-$6$ nontrivial $Z$-logical of the base code. We claim
$[u^\ast]\notin\img\cmap$.

Indeed, $u^\ast$ has nonzero \emph{seam flux}. Fix the seam $j=0$. For an appropriate
$X$-side dual cycle $\xi\in\ker H_X^{\mathsf T}$, the class functional
$\ell_\xi(w)=\xi^{\mathsf T}\partial_1^{\mathrm c}w$ of Corollary~\ref{cor:flux} evaluates on
$u^\ast$ to the parity of the seam-crossing incidences between $\xi$ and $\supp z^\ast$ --- a
finite count over the six cells of $\supp z^\ast$ --- and this parity is odd. Since $\ell_\xi$
vanishes on $\img\cmap$ (Corollary~\ref{cor:flux}, the easy inclusion
$\img\cmap\subseteq\ker\ell_\xi$), $\ell_\xi(u^\ast)\ne 0$ forces
$[u^\ast]\notin\img\cmap=\ker\tau_\ast$, where $\tau(u)=(u,u)$ is the transfer.

Since $[u^\ast]\notin\ker\tau_\ast$, the diagonal lift $\tau(u^\ast)=(u^\ast,u^\ast)$
is a nonzero class in $H_1(\mathrm{gross})$, i.e.\ a cover cycle that is not a
boundary. Its weight is $\wt{\tau(u^\ast)}=2\,\wt{u^\ast}=12$. Hence the gross code
carries a nontrivial $Z$-logical of weight exactly $12$, so $d(\mathrm{gross})\le 12$.

Combining the two bounds, $d(\mathrm{gross})=12$.
\end{proof}

Theorem~D closes the program: the base code's distance-$6$ geometry controls the
gross code's distance entirely, the dangerous sector through the factor-two floor of
Theorem~C and the safe sector through the Smith-class floor of
Proposition~\ref{thm:Mim}, with the diagonal lift of a single weight-$6$ base logical
showing the bound $12$ is achieved.

\end{document}
